% !TeX spellcheck = sv_FI
% !TeX encoding = UTF-8
\chapter{Antenner}
\label{sec:antenner}
\index{antenner}
\index{antenn}

Här ska vi diskutera några av de vanligaste förekommande antennerna
och särskilt från perspektivet att få till antenner som man kan bygga
själv på egen hand till sina radioapparater. Vi kommer främst titta på
antenner för kortvåg, VHF och lite UHF (låga delen, typiskt 70 cm).
En antenn man bygger själv är tillåten även med instegscertifikat.
Det som inte är tillåtet då är en hemmabyggd eller modifierad sändare,
se avsnitt~\ref{sec:insteg-restriktioner}.

\vfill
\minitoc
\clearpage

\section{Antenntyper}

Här beskriver vi några av de vanliga antenntyperna. Fokus är på
enkelhet och sådana man kan bygga själv även om också mer komplexa
antenner som Yagi-Uda omnämns.

\subsection{Dipolen}
\index{dipol}
\index{antenn!dipol}
\index{halvvågsdipol}
\index{antenn!höjd}
\index{NVIS!antenn}
\index{dipol!höjd}
\index{antenn!längd}
\index{halvvåg}

Dipolen är den antenn de flesta börjar med, och den som de andra
jämförs med. En halvvågstråd, matad på mitten. Den strålar mest rakt
ut från tråden, i ett åttamönster om man ser den ovanifrån. Längs
tråden är den tyst. En dipol för 40\,m som hänger öst--väst hörs
alltså bäst i riktningarna nord och syd. Men den strålar också rakt
uppåt och rakt nedåt.

Höjden bestämmer vart energin tar vägen. Hänger dipolen lågt, en bit
under en kvarts våglängd, går det mesta rakt upp, och det är
NVIS-läget, rätt när man vill nå runt om i landet och de närmaste
grannländerna. Kring en halv våglängd över marken lägger sig loben
flackare och räcker till DX. Blir den klart högre kan loben dela sig,
så att en del går flackt och en del brant.

\begin{table}[H]
\centering
\begin{tabular}{lrrl}
	\bf Band & \bf Kvartsvåg & \bf Halvvåg & \bf En lågt hängande tråd          \\ \hline
	160\,m   &         40\,m &       80\,m & nästan alltid NVIS                 \\
	80\,m    &         20\,m &       40\,m & NVIS hos de flesta                 \\
	60\,m    &         15\,m &       30\,m & NVIS, och det är ofta det man vill \\
	40\,m    &         10\,m &       20\,m & NVIS lågt, DX om den kommer upp    \\
	20\,m    &          5\,m &       10\,m & hygglig DX redan på 10--15\,m
\end{tabular}
\caption{Ungefärlig höjd över mark för NVIS och för ett flackare hopp.
Under kvartsvågen dominerar den branta loben.}
\label{tab:dipolhojd}
\end{table}

På 80\,m är en halv våglängd omkring 40\,m upp, så de flesta kör det
bandet lågt och får NVIS vare sig de tänkt på det eller inte. På 60\,m
är en kvarts våglängd omkring 15\,m, så samma tråd är en NVIS-antenn
där också. På 20\,m är en halv våglängd omkring 10\,m, och en dipol på
10--15\,m höjd är alltså NVIS på 80 och 60\,m och en hygglig DX-antenn
på 20\,m. Samma tråd, två olika jobb.

Den mer exakta längden står i teknikkapitlet, under antennlängd. En
tråddipol blir i praktiken kortare än den fria halvvågen, cirka
$142/f$ till $143/f$ meter när frekvensen är i MHz. Man klipper den
lite för lång och kortar den efter att man mätt stående våg. Ändarna
ska inte gå att nudda för människor eller djur eftersom där blir
spänningen hög när man sänder. Håll också tråden väl borta från
elledningar, undvik absolut att låta den gå över en väg eller gata och
tänk efter lite hur du placerar den.

    \begin{figure}[H]
    \centering
    \begin{tikzpicture}[
      x=1cm, y=1cm,
      antenn/.style={rhbrav, line width=1.35pt, line cap=round},
      koax/.style={black!80, line width=1.1pt, line cap=round},
      lina/.style={black!45, line width=0.6pt, dash pattern=on 2pt off 1.5pt},
      isolator/.style={fill=white, draw=black!75, line width=0.65pt, rounded corners=1pt},
      matt/.style={{Stealth[length=1.6mm]}-{Stealth[length=1.6mm]}, black!65, line width=0.55pt},
      nod/.style={fill=black, circle, inner sep=1.2pt}
    ]

      % ====================================================================
      % VÄNSTER DEL: KONSTRUKTION FRAMIFRÅN (MED STRÖM/SPÄNNING & HÖJD)
      % ====================================================================
      \begin{scope}[shift={(-2.35, 0)}]
        \node[etikett, font=\footnotesize\bfseries] at (0, 3.15) {Halvvågsdipol framifrån};

        % Mark och gräs/jord
        \fill[rhbjord] (-4.45, -2.35) rectangle (4.45, -1.95);
        \draw[marklinje] (-4.45, -1.95) -- (4.45, -1.95);

        % Två fästpunkter / stolpar i kanterna
        \draw[black!70, line width=1.4pt, line cap=round] (-4.15, -1.95) -- (-4.15, 1.55);
        \draw[black!70, line width=1.4pt, line cap=round] ( 4.15, -1.95) -- ( 4.15, 1.55);

        % Bärlinor mellan stolpar och ändisolatorer
        \draw[lina] (-4.15, 1.25) -- (-3.35, 1.25);
        \draw[lina] ( 3.35, 1.25) -- ( 4.15, 1.25);

        % Ändisolatorer
        \filldraw[isolator] (-3.35, 1.15) rectangle (-3.05, 1.35);
        \filldraw[isolator] ( 3.05, 1.15) rectangle ( 3.35, 1.35);
        \node[etikett, anchor=south] at (-3.20, 1.40) {Ändisolator};

        % Antennbenen (2 x lambda/4)
        \draw[antenn] (-3.05, 1.25) -- (-0.25, 1.25);
        \draw[antenn] ( 0.25, 1.25) -- ( 3.05, 1.25);
        \node[nod, fill=rhbrav] at (-0.25, 1.25) {};
        \node[nod, fill=rhbrav] at ( 0.25, 1.25) {};

        % Mittisolator + 1:1 strömbalun
        \filldraw[isolator] (-0.25, 1.12) rectangle (0.25, 1.38);
        \filldraw[fill=rhbjord!85, draw=black!75, line width=0.6pt, rounded corners=1pt]
          (-0.32, 0.52) rectangle (0.32, 0.98);
        \draw[black!75, line width=0.65pt] (-0.20, 1.25) -- (-0.15, 0.98);
        \draw[black!75, line width=0.65pt] ( 0.20, 1.25) -- ( 0.15, 0.98);
        \node[etikett, anchor=west] at (0.38, 0.75) {1:1-balun};

        % Koaxialkabel 50 ohm ner till radion
        \draw[koax] (0, 0.52) -- (0, -1.95);
        \node[etikett, anchor=west] at (0.12, -0.75) {Koax $50\,\Omega$};

        % Ström- och spänningsfördelning (ovanför tråden)
        % Ström I (blå streckad båge, max i mitten, noll i ändarna)
        \draw[rhbrav!75, line width=0.75pt, dash pattern=on 2.5pt off 1.5pt]
          (-3.05, 1.25) sin (0, 2.05) cos (3.05, 1.25);
        \node[etikett, text=rhbrav, anchor=south] at (0, 2.05) {Strömmax ($I$)};

        % Spänning U (rödbrun streckad kurva, max i ändarna, noll i mitten)
        \draw[rhbes!80, line width=0.75pt, dash pattern=on 1.5pt off 1.5pt]
          (-3.05, 2.00) cos (0, 1.25) sin (3.05, 0.50);
        \node[etikett, text=rhbes, anchor=south] at (2.35, 1.40) {Hög spänning i ändarna ($U$)};

        % Måttlinjer för benlängder (lambda/4) och totallängd (L ≈ 142,5 / f)
        \draw[matt] (-3.05, 0.15) -- (-0.05, 0.15)
          node[midway, fill=white, etikett] {$\approx \lambda/4$};
        \draw[matt] ( 0.05, 0.15) -- ( 3.05, 0.15)
          node[midway, fill=white, etikett] {$\approx \lambda/4$};

        \draw[matt] (-3.05, 2.60) -- (3.05, 2.60)
          node[midway, fill=white, etikett] {Total längd $L \approx \dfrac{142\text{--}143}{f\text{ [MHz]}}\text{ m}$ ($\lambda/2$)};
        \draw[black!40, line width=0.4pt] (-3.05, 1.40) -- (-3.05, 2.70);
        \draw[black!40, line width=0.4pt] ( 3.05, 1.40) -- ( 3.05, 2.70);

        % Höjd över mark (h)
        \draw[matt] (-3.75, -1.95) -- (-3.75, 1.25)
          node[midway, fill=white, etikett] {Höjd $h$};
      \end{scope}


      % ====================================================================
      % HÖGER DEL: STRÅLNINGSMÖNSTER OVANIFRÅN (ÅTTAMÖNSTER)
      % ====================================================================
      \begin{scope}[shift={(4.45, 0.10)}]
        \node[etikett, font=\footnotesize\bfseries] at (0, 3.05) {Strålning ovanifrån};

        % Kompasskors / hjälpcirkel
        \draw[black!20, line width=0.5pt] (0,0) circle[radius=1.85];
        \draw[black!30, line width=0.45pt, dash pattern=on 2pt off 2pt] (-2.10, 0) -- (2.10, 0);
        \draw[black!30, line width=0.45pt, dash pattern=on 2pt off 2pt] (0, -2.10) -- (0, 2.10);

        % Åttamönster (två lober vinkelrätt mot tråden: Nord och Syd)
        \filldraw[fill=rhbtack!22, draw=rhbtack!85!black, line width=1.05pt]
          (0,0)
          .. controls ( 0.95,  0.45) and ( 1.05,  1.85) .. (0,  1.85)
          .. controls (-1.05,  1.85) and (-0.95,  0.45) .. (0,0)
          .. controls ( 0.95, -0.45) and ( 1.05, -1.85) .. (0, -1.85)
          .. controls (-1.05, -1.85) and (-0.95, -0.45) .. (0,0);

        % Dipoltråden i mitten (öst-väst)
        \draw[antenn, line width=1.5pt] (-1.35, 0) -- (-0.10, 0);
        \draw[antenn, line width=1.5pt] ( 0.10, 0) -- ( 1.35, 0);
        \filldraw[isolator] (-0.10, -0.08) rectangle (0.10, 0.08);
        \fill[rhbrav] (-1.35, 0) circle[radius=0.045];
        \fill[rhbrav] ( 1.35, 0) circle[radius=0.045];

        % Etiketter för riktningar
        \node[etikett, anchor=south] at (0,  2.12) {\textbf{Nord} (max)};
        \node[etikett, anchor=north] at (0, -2.12) {\textbf{Syd} (max)};
        \node[etikett, anchor=west]  at ( 1.45, -0.22) {Tyst};
        \node[etikett, anchor=east]  at (-1.45, -0.22) {Tyst};
        \node[etikett, text=rhbrav, anchor=south] at (-0.75, 0.06) {Tråd};
      \end{scope}

    \end{tikzpicture}
    \caption{Halvvågsdipolen. Två lika långa ben om ungefär $\lambda/4$ matas i
    mitten via en strömbalun 1:1. I mitten är strömmen hög och impedansen nära
    $50\,\Omega$, medan spänningen blir hög i ändarna vid sändning. Ovanifrån
    strålar dipolen i ett åttamönster vinkelrätt mot tråden, och höjden $h$ över
    marken avgör om loben går brant uppåt (NVIS, $h < \lambda/4$) eller flackt
    (DX, $h \approx \lambda/2$).}
    \label{fig:dipol}
    \end{figure}



\subsubsection{Inverterad V}
\index{inverterad V}
\index{antenn!inverterad V}
\index{dipol!inverterad V}
\index{inverterad V!impedans}

Detta är egentligen samma dipol som ovan men placerad annorlunda.

Saknar man två fästpunkter räcker en. Mittpunkten hissas i en mast, ett
träd eller en nock, och benen går snett ner. Det är samma halvvåg, bara
vikt. Den tar mindre plats i sidled, och den tål att benen inte ligger
exakt i linje.

Impedansen sjunker när benen vinklas ner, och hamnar ofta närmare de
50\,$\Omega$ som koaxen vill se än vad den raka dipolen gör. Åttamönstret
finns kvar men är mjukare, så en inverterad V är snällare om man inte
kan välja riktning. Ändarna ska fortfarande upp från marken. Ligger de i
gräset försvinner effekt, och frekvensen drar iväg. En höjd där ingen
når dem räcker. Med ett spö eller ett träd i mitten är den en av de
antenner som är enklast att ta med ut.

\begin{figure}[H]
\centering
\begin{tikzpicture}[
  x=1cm, y=1cm,
  antenn/.style={rhbrav, line width=1.35pt, line cap=round},
  koax/.style={black!80, line width=1.1pt, line cap=round},
  lina/.style={black!45, line width=0.6pt, dash pattern=on 2pt off 1.5pt},
  isolator/.style={fill=white, draw=black!75, line width=0.65pt, rounded corners=1pt},
  matt/.style={{Stealth[length=1.6mm]}-{Stealth[length=1.6mm]}, black!65, line width=0.55pt},
  nod/.style={fill=black, circle, inner sep=1.2pt}
]

  % ====================================================================
  % VÄNSTER DEL: INVERTERAD V FRAMIFRÅN
  % ====================================================================
  \begin{scope}[shift={(-2.35, 0)}]
    \node[etikett, font=\footnotesize\bfseries] at (0, 3.05) {Inverterad V framifrån};

    % Mark och gräs/jord
    \fill[rhbjord] (-4.45, -2.35) rectangle (4.45, -1.95);
    \draw[marklinje] (-4.45, -1.95) -- (4.45, -1.95);

    % Enkel mittmast / spö i mitten
    \draw[black!65, line width=1.8pt, line cap=round] (0, -1.95) -- (0, 2.25);
    \node[etikett, anchor=east, text=black!75] at (-0.12, -1.05) {Mast / spö};

    % Låga fäststolpar / tältpinnar i sidorna
    \draw[black!70, line width=1.2pt, line cap=round] (-3.95, -1.95) -- (-3.95, -0.55);
    \draw[black!70, line width=1.2pt, line cap=round] ( 3.95, -1.95) -- ( 3.95, -0.55);

    % Mittisolator + 1:1 strömbalun i toppen
    \filldraw[isolator] (-0.25, 1.88) rectangle (0.25, 2.12);
    \filldraw[fill=rhbjord!85, draw=black!75, line width=0.6pt, rounded corners=1pt]
      (-0.30, 1.30) rectangle (0.30, 1.74);
    \draw[black!75, line width=0.65pt] (-0.18, 1.96) -- (-0.14, 1.74);
    \draw[black!75, line width=0.65pt] ( 0.18, 1.96) -- ( 0.14, 1.74);
    \node[etikett, anchor=west] at (0.36, 1.52) {1:1-balun ($\approx 50\,\Omega$)};

    % Antennbenen snett nedåt (ca 90–120 graders toppvinkel)
    \draw[antenn] (-0.25, 1.98) -- (-3.05, -0.35);
    \draw[antenn] ( 0.25, 1.98) -- ( 3.05, -0.35);
    \node[nod, fill=rhbrav] at (-0.25, 1.98) {};
    \node[nod, fill=rhbrav] at ( 0.25, 1.98) {};

    % Ändisolatorer
    \filldraw[isolator, rotate around={39.8:(-3.18,-0.46)}] (-3.33, -0.55) rectangle (-3.03, -0.37);
    \filldraw[isolator, rotate around={-39.8:(3.18,-0.46)}] ( 3.03, -0.55) rectangle ( 3.33, -0.37);
    \node[etikett, anchor=south east] at (-2.95, -0.30) {Ändisolator};

    % Bärlinor från ändisolatorer till fästpunkter
    \draw[lina] (-3.31, -0.57) -- (-3.95, -0.75);
    \draw[lina] ( 3.31, -0.57) -- ( 3.95, -0.75);

    % Koaxialkabel 50 ohm längs masten
    \draw[koax] (0.08, 1.30) -- (0.08, -1.95);
    \node[etikett, anchor=west] at (0.18, -0.25) {Koax $50\,\Omega$};

    % Toppvinkel (90–120 grader)
    \draw[black!60, line width=0.55pt] (-0.68, 1.62) arc[start angle=220, end angle=320, radius=0.88];
    \node[etikett, fill=white, inner sep=1pt] at (0, 0.88) {$90^\circ\text{--}120^\circ$};

    % Benlängder (lambda/4)
    \node[etikett, text=rhbrav, rotate=40, anchor=south] at (-1.65, 0.88) {$\approx \lambda/4$};
    \node[etikett, text=rhbrav, rotate=-40, anchor=south] at ( 1.65, 0.88) {$\approx \lambda/4$};

    % Ändhöjd över mark (fri från gräs och beröring)
    \draw[matt] (3.18, -1.95) -- (3.18, -0.55)
      node[midway, fill=white, etikett, anchor=east, xshift=-2pt] {Fri höjd $>2\text{ m}$};
  \end{scope}

  % ====================================================================
  % HÖGER DEL: MJUKARE STRÅLNINGSMÖNSTER OVANIFRÅN
  % ====================================================================
  \begin{scope}[shift={(4.45, 0.10)}]
    \node[etikett, font=\footnotesize\bfseries] at (0, 2.95) {Strålning ovanifrån};

    % Hjälpcirkel och axlar
    \draw[black!20, line width=0.5pt] (0,0) circle[radius=1.85];
    \draw[black!30, line width=0.45pt, dash pattern=on 2pt off 2pt] (-2.10, 0) -- (2.10, 0);
    \draw[black!30, line width=0.45pt, dash pattern=on 2pt off 2pt] (0, -2.10) -- (0, 2.10);

    % Rak dipol som referens (streckat åttamönster)
    \draw[black!40, line width=0.6pt, dash pattern=on 2pt off 1.5pt]
      (0,0)
      .. controls ( 0.95,  0.45) and ( 1.05,  1.85) .. (0,  1.85)
      .. controls (-1.05,  1.85) and (-0.95,  0.45) .. (0,0)
      .. controls ( 0.95, -0.45) and ( 1.05, -1.85) .. (0, -1.85)
      .. controls (-1.05, -1.85) and (-0.95, -0.45) .. (0,0);

    % Inverterad V: mjukare åttamönster som även fyller ut i sidled
    \filldraw[fill=rhbtack!22, draw=rhbtack!85!black, line width=1.05pt]
      (0, 1.72)
      .. controls ( 1.15,  1.72) and ( 1.05,  0.55) .. ( 0.68,  0.00)
      .. controls ( 1.05, -0.55) and ( 1.15, -1.72) .. ( 0.00, -1.72)
      .. controls (-1.15, -1.72) and (-1.05, -0.55) .. (-0.68,  0.00)
      .. controls (-1.05,  0.55) and (-1.15,  1.72) .. ( 0.00,  1.72);

    % Antennens projektion ovanifrån (kortare spann i sidled)
    \draw[antenn, line width=1.5pt] (-1.05, 0) -- (-0.10, 0);
    \draw[antenn, line width=1.5pt] ( 0.10, 0) -- ( 1.05, 0);
    \filldraw[isolator] (-0.10, -0.08) rectangle (0.10, 0.08);

    \node[etikett, anchor=south] at (0,  2.05) {Mjukare åtta};
    \node[etikett, anchor=north] at (0, -2.05) {(mindre djupa nollställen)};
  \end{scope}

\end{tikzpicture}
\caption{Inverterad V-dipol. En enda hög fästpunkt i mitten räcker och antennen
tar mindre plats i sidled än en rak dipol. När benen vinklas nedåt ($90^\circ\text{--}120^\circ$)
sjunker matningsimpedansen närmare $50\,\Omega$ och strålningsmönstret blir mer
rundstrålande. Ändarna ska hållas fritt över marken.}
\label{fig:inv-v}
\end{figure}

\subsubsection{Off-center-matad dipol}
\index{off-center}
\index{antenn!off-center}
\index{dipol!off-center}
\index{OCF}
\index{Guanella}
\index{strömbalun!Guanella}
\index{spänningsbalun}
\index{mantelström}
\index{common mode}
\index{Ruthroff}
\index{balun!4:1}
\index{balun!Guanella}
\index{OCF!benlängd}
\index{4:1}

Den off-center-matade dipolen är vägen när man vill ha flera band på
koax. På engelska heter den \textit{off-center fed dipole}, OCF. Det
är samma halvvågstråd som den vanliga dipolen, men matningen sitter en
tredjedel in från ena änden. Korta benet är en tredjedel av tråden,
långa benet två tredjedelar.

Hela längden $L$ är tråddipolens, $142/f$ till $143/f$ meter när
frekvensen är i MHz, samma tal som i teknikkapitlet. Benen blir

\begin{align}
	l_{\text{kort}} &= \frac{L}{3} \\
	l_{\text{lång}} &= \frac{2L}{3}
\end{align}

På 3,6\,MHz, mitt på 80\,m, är tråden knappt 40\,m om man räknar
$143/f$. Korta benet blir då drygt 13\,m och långa benet drygt 26\,m.
Man klipper lite långt och kortar efter att stående våg är mätt,
precis som på en vanlig dipol. Klipper man den för 40\,m blir det en
mindre antenn för 40, 20 och 10\,m.

En mittmatad dipol har låg impedans på grundtonen och på de udda
övertonerna. På jämna övertoner ligger mitten i ett spänningsmaximum,
och impedansen är flera tusen ohm. Vid en tredjedels längd ligger
impedansen kring 200\,$\Omega$ både på grundtonen och på de jämna
övertonerna. En omsättning 4:1 tar ner det mot 50\,$\Omega$. Koaxen
kan då gå hela vägen till radion med måttlig stående våg. En OCF för
80\,m går därför bra på 80, 40, 20 och 10\,m. 15\,m och WARC-banden är
ojämnare och vill ofta ha en tuner. 4:1 ska sitta uppe i
matningspunkten. Sitter den nere vid radion är det koaxen som bär de
200\,$\Omega$, och då är förlusten tillbaka.

Effektiviteten kommer av att hela tråden strålar. Inga spolar och inga
traps som stjäl effekt och gör bandet smalt. På de band den är gjord
för är den i samma klass som en vanlig dipol. En enda tråd räcker till
flera av dem, och koaxen värmer inte bort det man just vunnit.

Det som avgör om det blir så i praktiken är balunen. Antennen är
osymmetrisk. Benen är olika långa, så strömmarna vill inte vara lika.
En spänningsbalun, oftast en Ruthroff 4:1, tvingar fram lika
spänningar och har sämre undertryckning av common mode. Matarledningen
blir då en del av det som strålar och loben drar snett, bruset ökar
och man får radiofrekvens i rummet. En del av effekten värmer kärnan i
stället för att gå ut i tråden.

En strömbalun av Guanella-typ håller emot det. Den är byggd av två
transmissionsledningar, kopplade så att de ger 4:1, och de släpper i
praktiken bara igenom ström som går åt motsatt håll i de två
ledarna. Common mode-ström, ström åt samma håll i både skärmen och
mittenledaren, möter där ferritens drossel som ett växelströmsmotstånd
och undertryckningen är klart bättre än i en spänningsbalun. Därför är
det helst en Guanella man sätter i matningspunkten. OCF är så pass
sned att en ensam 4:1 inte alltid räcker. En extra strömbalun 1:1 på
koaxen strax under den hjälper också för att minska common
mode-problematiken.

Höjden följer den vanliga dipolen. Lågt blir det NVIS, högre blir
loben flackare. Åttamönstret är lite snett eftersom benen är olika
långa, och på de högre banden delar loben sig.

Förutom obalansen i antennen är den fantastisk när man enkelt behöver
en antenn som snabbt kan köras på flera band. Den är dock inte
särskilt kompakt men den ger ofta mycket bra resultat och den är lätt
att bygga.

\begin{figure}[H]
\centering
\begin{tikzpicture}[
  x=1cm, y=1cm,
  antenn/.style={rhbrav, line width=1.35pt, line cap=round},
  koax/.style={black!80, line width=1.1pt, line cap=round},
  lina/.style={black!45, line width=0.6pt, dash pattern=on 2pt off 1.5pt},
  isolator/.style={fill=white, draw=black!75, line width=0.65pt, rounded corners=1pt},
  matt/.style={{Stealth[length=1.6mm]}-{Stealth[length=1.6mm]}, black!65, line width=0.55pt},
  nod/.style={fill=black, circle, inner sep=1.2pt}
]

  % ====================================================================
  % VÄNSTER DEL: OCF-DIPOLENS UPPBYGGNAD (1/3 OCH 2/3)
  % ====================================================================
  \begin{scope}[shift={(-2.25, 0)}]
    \node[etikett, font=\footnotesize\bfseries] at (0, 3.10) {Off-center-matad dipol (OCF)};

    % Mark och gräs/jord
    \fill[rhbjord] (-4.55, -2.35) rectangle (4.55, -1.95);
    \draw[marklinje] (-4.55, -1.95) -- (4.55, -1.95);

    % Två fästpunkter / stolpar i kanterna
    \draw[black!70, line width=1.4pt, line cap=round] (-4.25, -1.95) -- (-4.25, 1.65);
    \draw[black!70, line width=1.4pt, line cap=round] ( 4.25, -1.95) -- ( 4.25, 1.65);

    % Bärlinor och ändisolatorer
    \draw[lina] (-4.25, 1.35) -- (-3.60, 1.35);
    \draw[lina] ( 3.60, 1.35) -- ( 4.25, 1.35);
    \filldraw[isolator] (-3.60, 1.25) rectangle (-3.30, 1.45);
    \filldraw[isolator] ( 3.30, 1.25) rectangle ( 3.60, 1.45);

    % Matningspunkt vid 1/3 av längden (x = -1.10, total tråd -3.30 till +3.30 = 6.60)
    % Kort ben: -3.30 till -1.32 (längd ~2.20), Långt ben: -0.88 till +3.30 (längd ~4.40)
    \draw[antenn] (-3.30, 1.35) -- (-1.32, 1.35);
    \draw[antenn] (-0.88, 1.35) -- ( 3.30, 1.35);
    \node[nod, fill=rhbrav] at (-1.32, 1.35) {};
    \node[nod, fill=rhbrav] at (-0.88, 1.35) {};

    % Isolator vid matningspunkten (x = -1.10)
    \filldraw[isolator] (-1.32, 1.22) rectangle (-0.88, 1.48);
    \node[etikett, anchor=south, text=black!75] at (-1.10, 1.52) {$\approx 200\,\Omega$};

    % Guanella 4:1 strömbalun direkt i matningspunkten
    \filldraw[fill=rhbjord!85, draw=black!75, line width=0.6pt, rounded corners=1pt]
      (-1.45, 0.58) rectangle (-0.75, 1.05);
    \draw[black!75, line width=0.65pt] (-1.28, 1.35) -- (-1.22, 1.05);
    \draw[black!75, line width=0.65pt] (-0.92, 1.35) -- (-0.98, 1.05);
    \node[etikett, anchor=west] at (-0.68, 0.82) {Guanella 4:1 ($200 \to 50\,\Omega$)};

    % Koax 50 ohm nedåt + extra 1:1 strömbalun (mantelströmsdrossel)
    \draw[koax] (-1.10, 0.58) -- (-1.10, -0.02);
    \filldraw[fill=rhbe!45, draw=black!75, line width=0.6pt, rounded corners=1pt]
      (-1.40, -0.45) rectangle (-0.80, -0.02);
    \node[etikett, anchor=west] at (-0.68, -0.24) {Extra 1:1-balun (mot mantelström)};
    \draw[koax] (-1.10, -0.45) -- (-1.10, -1.95);
    \node[etikett, anchor=west] at (-0.98, -1.25) {Koax $50\,\Omega$};

    % Måttlinjer för korta benet (L/3) och långa benet (2L/3)
    \draw[matt] (-3.30, 2.02) -- (-1.15, 2.02)
      node[midway, fill=white, etikett] {$l_{\text{kort}} = \dfrac{L}{3}$};
    \draw[matt] (-1.05, 2.02) -- ( 3.30, 2.02)
      node[midway, fill=white, etikett] {$l_{\text{lång}} = \dfrac{2L}{3}$};

    % Totallängd L
    \draw[matt] (-3.30, 2.62) -- ( 3.30, 2.62)
      node[midway, fill=white, etikett] {Total längd $L \approx \dfrac{142\text{--}143}{f\text{ [MHz]}}\text{ m}$};
    \draw[black!40, line width=0.4pt] (-3.30, 1.50) -- (-3.30, 2.72);
    \draw[black!40, line width=0.4pt] (-1.10, 1.78) -- (-1.10, 2.12);
    \draw[black!40, line width=0.4pt] ( 3.30, 1.50) -- ( 3.30, 2.72);
  \end{scope}

  % ====================================================================
  % HÖGER DEL: IMPEDANS VID 1/3 PÅ GRUNDTON OCH ÖVERTONER
  % ====================================================================
  \begin{scope}[shift={(4.55, 0.05)}]
    \node[etikett, font=\footnotesize\bfseries] at (0, 3.05) {Strömfördelning vid $1/3$};

    % Vertikal linje för 1/3-matningspunkten
    \draw[rhbes, line width=0.7pt, dash pattern=on 2.5pt off 1.5pt] (-0.60, -1.65) -- (-0.60, 2.35);
    \node[etikett, text=rhbes, anchor=south] at (-0.60, 2.35) {Matning ($1/3$)};

    % 1) Grundton (80 m / 40 m): halvvåg
    \draw[black!55, line width=0.8pt] (-1.80, 1.45) -- (1.80, 1.45);
    \draw[rhbrav, line width=0.95pt] (-1.80, 1.45) sin (0, 2.10) cos (1.80, 1.45);
    \fill[rhbes] (-0.60, 2.01) circle[radius=0.055];
    \node[etikett, anchor=north west] at (-1.80, 1.42) {$f_0$ ($\lambda/2$): $\approx 200\,\Omega$};

    % 2) Andra deltonen (40 m / 20 m): helvåg (mitten har I=0, men 1/3 har god ström!)
    \draw[black!55, line width=0.8pt] (-1.80, 0.10) -- (1.80, 0.10);
    \draw[rhbrav, line width=0.95pt]
      (-1.80, 0.10) sin (-0.90, 0.72) cos (0, 0.10) sin (0.90, -0.52) cos (1.80, 0.10);
    \fill[rhbes] (-0.60, 0.64) circle[radius=0.055];
    \node[etikett, anchor=north west] at (-1.80, -0.52) {$2f_0$ ($1\lambda$): $\approx 200\,\Omega$};

    % 3) Fjärde deltonen (20 m / 10 m): 2 våglängder
    \draw[black!55, line width=0.8pt] (-1.80, -1.35) -- (1.80, -1.35);
    \draw[rhbrav, line width=0.95pt]
      (-1.80, -1.35) sin (-1.35, -0.80) cos (-0.90, -1.35)
      sin (-0.45, -1.90) cos ( 0.00, -1.35)
      sin ( 0.45, -0.80) cos ( 0.90, -1.35)
      sin ( 1.35, -1.90) cos ( 1.80, -1.35);
    \fill[rhbes] (-0.60, -1.83) circle[radius=0.055];
    \node[etikett, anchor=north west] at (-1.80, -1.92) {$4f_0$ ($2\lambda$): $\approx 200\,\Omega$};
  \end{scope}

\end{tikzpicture}
\caption{Off-center-matad dipol (OCF). Matningspunkten sitter en tredjedel in
från ena änden ($l_{\text{kort}} = L/3$, $l_{\text{lång}} = 2L/3$). Där ligger
impedansen kring $200\,\Omega$ både på grundfrekvensen $f_0$ och på de jämna
övertonerna ($2f_0$, $4f_0$). En Guanella 4:1-strömbalun i matningspunkten
transformerar ner till $50\,\Omega$, och en extra 1:1-strömbalun strax under
stoppar mantelströmmar på koaxialkabeln.}
\label{fig:ocf}
\end{figure}

\subsubsection{Stegmatad dipol, doublet}
\index{doublet}
\index{antenn!doublet}
\index{stegmatning}
\index{antenn!stegmatning}
\index{stege!doublet}
\index{stege}
\index{doublet!längd}
\index{antenntuner!doublet}
\index{strömbalun!doublet}

En doublet är en dipol som inte är klippt för ett enda band. Den matas
på mitten med öppen stege, kring 450 eller 600\,$\Omega$, hela vägen
in till en balanserad avstämningsenhet. Koax hör inte hemma på den
sträckan. Har man bara en vanlig tuner sätter man en strömbalun 1:1
alldeles intill den, och en kort koaxstump därifrån till radion.

Det som gör stegen värd besväret är förlusten. Den är liten även när
stående våg är hög. Samma missanpassning i en koax blir värme i
kabeln, och på de högre banden kan en god del av effekten försvinna
innan den når tråden. I stegen kommer den fram. Tunern vid radion får
slutsteget att se 50\,$\Omega$. Missanpassningen får stanna i stegen,
där den kostar lite. Håll stegen fri från metall. Lindas den längs en
mast eller tejpas mot en ränna börjar den uppföra sig som en koax, och
då är den låga förlusten borta.

Därför ska tråden ligga mellan resonanserna på de band man vill köra.
Inte på en halvvåg, inte på en helvåg, och inte på en multipel av dem.
Vid halvvåg och vid udda multiplar, tre halvvågor och så vidare,
ligger matningen i ett strömmaximum. Impedansen blir låg, ofta under
hundra $\Omega$, mot en stege som är gjord för flera hundra. Vid
helvåg och vid jämna multiplar ligger matningen i ett
spänningsmaximum.  Impedansen går då upp i flera tusen $\Omega$. Båda
lägena är svåra för tunern. Spänningen eller strömmen i stegen blir
onödigt hög, och en del av det man vann på den snåla stegen försvinner
i anpassningen.

Längden i tabellen är hela tråden, från ände till ände, matad på
mitten. Den är räknad med samma förkortning som tråddipolen i
teknikkapitlet, omkring $142/f$ till $143/f$, och längderna ligger
också fria om man räknar den fria halvvågen $150/f$. Halvvågen på
lägsta bandet, den man annars klipper en vanlig dipol efter, är just
den längd man lämnar här. Bandgränserna är de svenska, och WARC-banden
är med. Kortare än omkring en tredjedels våg på lägsta bandet blir
tråden klen. Då hjälper inte stegen, för mer av effekten fastnar i
tunern.

En halvmeters skillnad kan räcka för att hamna på en resonans,
särskilt på 10\,m där bandet är brett. 27,5\,m ska hållas nära det
talet.  28\,m landar på 60\,m. 32,5\,m tål ett par
decimeter. Längderna för 160\,m tål ungefär en meter, och 13 samt
18,5\,m tål några decimeter.

\begin{table}[H]
\centering
\begin{tabular}{lll}
	\textbf{Band} & \textbf{Längd} & \textbf{Läget}                                   \\ \hline
	80--10\,m     & 27,5\,m        & Kort. Drygt en tredjedels våg på 80\,m.          \\
	80--10\,m     & 32,5\,m        & Lagom. Klart under en halv våg på 80\,m.         \\
	160--20\,m    & 48\,m          & Kort, när tråden inte får bli lång.              \\
	160--20\,m    & 58\,m          & Lagom på 160\,m.                                 \\
	160--20\,m    & 66\,m          & Längre, fortfarande under en halv våg på 160\,m. \\
	40--6\,m      & 13\,m          & Kort. Omkring en tredjedels våg på 40\,m.        \\
	40--6\,m      & 18,5\,m        & En bit under en halv våg på 40\,m.               \\
\end{tabular}
\caption{Lämpliga längder för stegmatad doublet}
\end{table}

Höjden gör samma sak som för en vanlig dipol. Lågt ger NVIS på de låga
banden, högre ger flackare lob och DX. På de högre banden blir tråden
flera våglängder och loben delar sig. Det är normalt för en doublet.

\begin{figure}[H]
\centering
\begin{tikzpicture}[
  x=1cm, y=1cm,
  antenn/.style={rhbrav, line width=1.35pt, line cap=round},
  stege/.style={rhbrav!90!black, line width=0.9pt, line cap=round},
  koax/.style={black!80, line width=1.1pt, line cap=round},
  lina/.style={black!45, line width=0.6pt, dash pattern=on 2pt off 1.5pt},
  isolator/.style={fill=white, draw=black!75, line width=0.65pt, rounded corners=1pt},
  matt/.style={{Stealth[length=1.6mm]}-{Stealth[length=1.6mm]}, black!65, line width=0.55pt},
  nod/.style={fill=black, circle, inner sep=1.2pt}
]

  % ====================================================================
  % VÄNSTER DEL: STEGMATAD DOUBLET FRAMIFRÅN
  % ====================================================================
  \begin{scope}[shift={(-2.30, 0)}]
    \node[etikett, font=\footnotesize\bfseries] at (0, 3.10) {Stegmatad doublet med tuner};

    % Mark och gräs/jord
    \fill[rhbjord] (-4.50, -2.45) rectangle (4.50, -2.05);
    \draw[marklinje] (-4.50, -2.05) -- (4.50, -2.05);

    % Två fästpunkter / stolpar i kanterna
    \draw[black!70, line width=1.4pt, line cap=round] (-4.20, -2.05) -- (-4.20, 1.75);
    \draw[black!70, line width=1.4pt, line cap=round] ( 4.20, -2.05) -- ( 4.20, 1.75);

    % Bärlinor och ändisolatorer
    \draw[lina] (-4.20, 1.45) -- (-3.45, 1.45);
    \draw[lina] ( 3.45, 1.45) -- ( 4.20, 1.45);
    \filldraw[isolator] (-3.45, 1.35) rectangle (-3.15, 1.55);
    \filldraw[isolator] ( 3.15, 1.35) rectangle ( 3.45, 1.55);

    % Två lika långa antennben (icke-resonanta, L/2 vardera)
    \draw[antenn] (-3.15, 1.45) -- (-0.22, 1.45);
    \draw[antenn] ( 0.22, 1.45) -- ( 3.15, 1.45);
    \node[nod, fill=rhbrav] at (-0.22, 1.45) {};
    \node[nod, fill=rhbrav] at ( 0.22, 1.45) {};

    % Mittisolator (ingen balun i toppen!)
    \filldraw[isolator] (-0.22, 1.32) rectangle (0.22, 1.58);

    % Öppen stege (450–600 ohm) rakt ner till tunern
    \draw[stege] (-0.16, 1.45) -- (-0.16, -0.78);
    \draw[stege] ( 0.16, 1.45) -- ( 0.16, -0.78);
    \foreach \y in {1.15, 0.80, 0.45, 0.10, -0.25, -0.60} {
      \draw[black!65, line width=0.65pt] (-0.18, \y) -- (0.18, \y);
    }
    \node[etikett, anchor=west] at (0.28, 0.55) {Öppen stege $450\text{--}600\,\Omega$};
    \node[etikett, anchor=west, text=black!70] at (0.28, 0.18) {(låg förlust trots hög SWR,};
    \node[etikett, anchor=west, text=black!70] at (0.28, -0.15) {hålls fri från metall)};

    % 1:1 strömbalun direkt vid tunern + antenntuner + kort koax till radio
    \filldraw[fill=rhbjord!85, draw=black!75, line width=0.6pt, rounded corners=1pt]
      (-0.34, -1.25) rectangle (0.34, -0.78);
    \node[etikett, anchor=east] at (-0.42, -1.01) {1:1-balun};

    % Antenntuner (ATU)
    \filldraw[fill=rhbfsvag!75, draw=black!75, line width=0.65pt, rounded corners=1.5pt]
      (-0.75, -1.90) rectangle (0.75, -1.38);
    \draw[koax] (0, -1.25) -- (0, -1.38);
    \node[etikett, font=\scriptsize\bfseries] at (0, -1.64) {Antenntuner};

    % Kort koax 50 ohm till sändtagare (TRX)
    \draw[koax] (0.75, -1.64) -- (1.85, -1.64);
    \node[etikett, anchor=south] at (1.30, -1.58) {Kort koax};
    \node[etikett, anchor=north] at (1.30, -1.68) {$50\,\Omega$};

    \filldraw[fill=white, draw=black!75, line width=0.65pt, rounded corners=1.5pt]
      (1.85, -1.90) rectangle (3.10, -1.38);
    \node[etikett, font=\scriptsize\bfseries] at (2.475, -1.64) {Radio};

    % Måttlinjer för benen och totallängd (vald utanför halvvåg/helvåg)
    \draw[matt] (-3.15, 2.02) -- (-0.05, 2.02)
      node[midway, fill=white, etikett] {$L/2$};
    \draw[matt] ( 0.05, 2.02) -- ( 3.15, 2.02)
      node[midway, fill=white, etikett] {$L/2$};

    \draw[matt] (-3.15, 2.60) -- ( 3.15, 2.60)
      node[midway, fill=white, etikett] {Icke-resonant längd $L$ (t.ex.\ $27{,}5\text{ m}$ eller $32{,}5\text{ m}$)};
    \draw[black!40, line width=0.4pt] (-3.15, 1.60) -- (-3.15, 2.70);
    \draw[black!40, line width=0.4pt] ( 3.15, 1.60) -- ( 3.15, 2.70);
  \end{scope}

  % ====================================================================
  % HÖGER DEL: STRÅLNINGSMÖNSTER PÅ LÅGA RESPEKTIVE HÖGA BAND
  % ====================================================================
  \begin{scope}[shift={(4.50, 0.05)}]
    \node[etikett, font=\footnotesize\bfseries] at (0, 3.05) {Lober ovanifrån};

    % Hjälpcirkel och axlar
    \draw[black!20, line width=0.5pt] (0,0) circle[radius=1.80];
    \draw[black!30, line width=0.45pt, dash pattern=on 2pt off 2pt] (-2.05, 0) -- (2.05, 0);
    \draw[black!30, line width=0.45pt, dash pattern=on 2pt off 2pt] (0, -2.05) -- (0, 2.05);

    % Låga band (streckad grön åtta: som vanlig dipol)
    \draw[rhbtack!85!black, line width=0.85pt, dash pattern=on 3pt off 1.8pt]
      (0,0)
      .. controls ( 0.90,  0.42) and ( 1.00,  1.78) .. (0,  1.78)
      .. controls (-1.00,  1.78) and (-0.90,  0.42) .. (0,0)
      .. controls ( 0.90, -0.42) and ( 1.00, -1.78) .. (0, -1.78)
      .. controls (-1.00, -1.78) and (-0.90, -0.42) .. (0,0);

    % Högre band (L > lambda): loben delar sig i fyra huvudlober (klöverblad)
    \filldraw[fill=rhbf!25, draw=rhbrav, line width=1.0pt]
      (0,0)
      .. controls ( 0.35,  1.05) and ( 1.15,  1.65) .. ( 1.32,  1.32)
      .. controls ( 1.65,  1.15) and ( 1.05,  0.35) .. (0,0)
      .. controls ( 1.05, -0.35) and ( 1.65, -1.15) .. ( 1.32, -1.32)
      .. controls ( 1.15, -1.65) and ( 0.35, -1.05) .. (0,0)
      .. controls (-0.35, -1.05) and (-1.15, -1.65) .. (-1.32, -1.32)
      .. controls (-1.65, -1.15) and (-1.05, -0.35) .. (0,0)
      .. controls (-1.05,  0.35) and (-1.65,  1.15) .. (-1.32,  1.32)
      .. controls (-1.15,  1.65) and (-0.35,  1.05) .. (0,0);

    % Dipoltråden i mitten
    \draw[antenn, line width=1.5pt] (-1.35, 0) -- (-0.10, 0);
    \draw[antenn, line width=1.5pt] ( 0.10, 0) -- ( 1.35, 0);
    \filldraw[isolator] (-0.10, -0.08) rectangle (0.10, 0.08);

    \node[etikett, text=rhbtack!85!black, anchor=south] at (0, 2.02) {Låga band: åtta};
    \node[etikett, text=rhbrav, anchor=north] at (0, -2.02) {Höga band: delad lob};
  \end{scope}

\end{tikzpicture}
\caption{Stegmatad doublet för flerbandsdrift. Trådlängden $L$ väljs så att den
undviker exakta halv- och helvågsresonanser. Den öppna stegen ($450\text{--}600\,\Omega$)
har mycket låga förluster även vid hög stående våg och går hela vägen ner till
en balanserad antenntuner (eller en 1:1-strömbalun direkt vid tunern). På högre
band där tråden är längre än en våglängd delar sig strålningsloben.}
\label{fig:doublet}
\end{figure}

\subsection{Vertikalantenn och jordplan}
\index{vertikalantenn}
\index{antenn!vertikal}
\index{jordplan}
\index{antenn!jordplan}
\index{radial}
\index{kvartsvåg}
\index{vertikalantenn!radial}
\index{vertikalantenn!portabel}

En kvartsvåg som står upp strålar runt om, och lågt över horisonten
när den har något att arbeta mot. Det låga skottet är det man vill ha
för DX på de högre kortvågsbanden. Längden är cirka $71/f$ meter,
samma tumregel som i teknikkapitlet.

Den behöver ett jordplan. Står den på marken är radialerna det som
avgör hur mycket av effekten som blir radio i stället för värme i
backen. På marken behöver de inte vara avstämda. De bildar ett nät,
och fler är bättre. En bit upp i luften ska de vara ungefär en
kvartsvåg, tre eller fyra stycken, gärna sluttande nedåt. Då hamnar
matningen ofta nära 50\,$\Omega$.

Vertikalen är som regel brusigare på mottagning än en horisontell
dipol. Mycket av bruset runt husen är vertikalt polariserat. Över
salt vatten, som i avsnittet om markvåg, är den däremot
svårslagen. Där finns jordplanet redan. På ett öppet fält, med tre
eller fyra radialer och ett spö, är den ett bra val för portabelt när
en dipol inte får plats.

En enklare sådan kan konstrueras av ett fiskespö. Den matade tråden
ska vara ungefär 1/4 av våglängden (beräkna som ena benet på en
halvvåg). På marken lägger man sedan ut en motvikt med ungefär samma
längd i den riktning man vill ha bästa verkan av antennen, eller så
har man flera motvikter, exempelvis 3 stycken arrangerade som ett Y
och får den då ganska rundstrålande. Motvikterna kan vara kortare än
det vertikala men är de ungefär lika långa får man bästa effekten.

\begin{figure}[H]
\centering
\begin{tikzpicture}[
  x=1cm, y=1cm,
  antenn/.style={rhbrav, line width=1.4pt, line cap=round},
  radial/.style={rhbtack!85!black, line width=1.15pt, line cap=round},
  koax/.style={black!80, line width=1.1pt, line cap=round},
  isolator/.style={fill=white, draw=black!75, line width=0.65pt, rounded corners=1pt},
  matt/.style={{Stealth[length=1.6mm]}-{Stealth[length=1.6mm]}, black!65, line width=0.55pt},
  nod/.style={fill=black, circle, inner sep=1.2pt}
]

  % ====================================================================
  % VÄNSTER DEL: MARKMONTERAD / FISKESPÖ MED MARKRADIALER + LÅGA DX-LOBER
  % ====================================================================
  \begin{scope}[shift={(-3.35, 0)}]
    \node[etikett, font=\footnotesize\bfseries] at (0, 3.15) {På marken (t.ex.\ fiskespö)};

    % Mark och jord
    \fill[rhbjord] (-3.25, -2.35) rectangle (3.25, -1.85);
    \draw[marklinje] (-3.25, -1.85) -- (3.25, -1.85);

    % Låga strålningslober (flack vinkel för DX åt alla håll)
    \filldraw[fill=rhbtack!18, draw=rhbtack!80!black, line width=0.8pt, dash pattern=on 2.5pt off 1.5pt]
      (0, -1.75)
      .. controls ( 0.85, -1.65) and ( 2.85, -1.35) .. ( 3.05, -0.55)
      .. controls ( 3.15,  0.15) and ( 1.45, -0.45) .. (0, -1.75);
    \filldraw[fill=rhbtack!18, draw=rhbtack!80!black, line width=0.8pt, dash pattern=on 2.5pt off 1.5pt]
      (0, -1.75)
      .. controls (-0.85, -1.65) and (-2.85, -1.35) .. (-3.05, -0.55)
      .. controls (-3.15,  0.15) and (-1.45, -0.45) .. (0, -1.75);
    \node[etikett, text=rhbtack!80!black, anchor=south] at (2.15, -0.42) {Låg lob (DX)};

    % Fiskespö / glasfibermast som bär tråden
    \draw[black!35, line width=1.8pt, line cap=round] (-0.08, -1.85) -- (-0.08, 2.35);
    \node[etikett, text=black!65, anchor=east] at (-0.18, 1.55) {Spö};

    % Vertikal strålare (lambda/4 ≈ 71 / f)
    \draw[antenn] (0.04, -1.65) -- (0.04, 2.25);
    \node[nod, fill=rhbrav] at (0.04, -1.65) {};

    % Basisolator / matningspunkt
    \filldraw[isolator] (-0.18, -1.85) rectangle (0.22, -1.65);

    % Radialer / motvikter utlagda på marken (Y-form / nät)
    \draw[radial] (-0.18, -1.82) -- (-2.75, -1.82);
    \draw[radial] ( 0.22, -1.82) -- ( 2.75, -1.82);
    \draw[radial] (-0.05, -1.85) -- (-1.55, -2.25);
    \draw[radial] ( 0.10, -1.85) -- ( 1.65, -2.25);
    \node[etikett, text=rhbtack!85!black, anchor=north] at (0, -2.35)
      {Radialer / motvikter på marken (t.ex.\ 3 st i Y eller radialnät)};

    % Koax 50 ohm
    \draw[koax] (0.22, -1.75) -- (1.45, -1.75);
    \node[etikett, anchor=south] at (1.05, -1.72) {Koax $50\,\Omega$};

    % Måttlinje för kvartsvågen
    \draw[matt] (-0.85, -1.65) -- (-0.85, 2.25)
      node[midway, fill=white, etikett] {$\approx \dfrac{\lambda}{4} \approx \dfrac{71}{f}\text{ m}$};
    \draw[black!40, line width=0.4pt] (-0.95, 2.25) -- (0.00, 2.25);
  \end{scope}

  % ====================================================================
  % HÖGER DEL: UPPHÖJD VERTIKAL MED SLUTTANDE AVSTÄMDA RADIALER (50 OHM)
  % ====================================================================
  \begin{scope}[shift={(3.55, 0)}]
    \node[etikett, font=\footnotesize\bfseries] at (0, 3.15) {Upphöjd med jordplan};

    % Mark och jord
    \fill[rhbjord] (-3.05, -2.35) rectangle (3.05, -1.85);
    \draw[marklinje] (-3.05, -1.85) -- (3.05, -1.85);

    % Bärrör / mast upp i luften
    \draw[black!65, line width=1.6pt, line cap=round] (0, -1.85) -- (0, -0.35);

    % Basisolator i matningspunkten
    \filldraw[isolator] (-0.20, -0.35) rectangle (0.20, -0.12);

    % Vertikal strålare (lambda/4)
    \draw[antenn] (0, -0.12) -- (0, 2.45);
    \node[nod, fill=rhbrav] at (0, -0.12) {};
    \node[etikett, text=rhbrav, anchor=west] at (0.12, 1.35) {Strålare $\approx \lambda/4$};

    % 3–4 sluttande avstämda radialer (ger ca 50 ohm)
    \draw[radial] (-0.20, -0.25) -- (-2.25, -1.25);
    \draw[radial] ( 0.20, -0.25) -- ( 2.25, -1.25);
    \draw[radial, line width=0.95pt] (-0.08, -0.30) -- (-0.95, -1.42);
    \draw[radial, line width=0.95pt] ( 0.08, -0.30) -- ( 1.05, -1.38);

    \node[etikett, text=rhbtack!85!black, rotate=26, anchor=south] at (-1.35, -0.75) {$\approx \lambda/4$};
    \node[etikett, text=rhbtack!85!black, rotate=-26, anchor=south] at ( 1.35, -0.75) {$\approx \lambda/4$};

    % Koax 50 ohm längs masten
    \draw[koax] (0.08, -0.35) -- (0.08, -1.85);
    \node[etikett, anchor=west] at (0.18, -1.55) {Koax $50\,\Omega$};

    \node[etikett, anchor=north] at (0, -2.35)
      {3--4 sluttande $\lambda/4$-radialer $\Rightarrow Z \approx 50\,\Omega$};
  \end{scope}

\end{tikzpicture}
\caption{Kvartsvågs vertikalantenn ($\approx 71/f\text{ m}$). Antennen är
rundstrålande i horisontalplanet med låg strålningsvinkel för DX och kräver ett
jordplan att arbeta mot. Till vänster: markmonterad vertikal (t.ex.\ tråd längs
ett fiskespö) med radialer/motvikter på marken. Till höger: upphöjd vertikal där
3--4 sluttande kvartsvågsradialer ger en matningsimpedans nära $50\,\Omega$.}
\label{fig:vertikal}
\end{figure}

\subsection{Longwire och ändmatad, EFHW och EFRW}
\index{longwire}
\index{EFRW}
\index{antenn!longwire}
\index{endmatad}
\index{antenn!endmatad}
\index{EFHW}
\index{endmatad!halvvåg}
\index{motvikt}
\index{överton}
\index{49:1}
\index{transformator!endmatad}

Vi börjar med att reda ut några begrepp. Det finns i huvudsak två
sorters ändmatade antenner med två olika filosofier och dessa är EFHW
som är en resonant antenn med mycket hög matningsimpedans, kräver en
unun med högt omsättningstal (49:1 är vanligt) och fungerar på band
som är resonanta mot halvvågen samt helvågen för dess längd.

Den andra sorten är det som ibland kallas ``random wire'' men den är
inte så random, tvärt om är det rätt viktigt att klippa till den så
den inte kommer i resonans. Den får lägre matningsimpedans, klarar sig
med en unun på runt 9:1 i omsättningstal och kan med fördel köras med
tuner och får relativt goda egenskaper. Den fungerar bra på alla
frekvenser som den INTE är resonant på.

En tråd som är flera våglängder lång strålar i lober längs tråden och
inte tvärs den som den vanliga dipolen när man kör den på frekvenser
som ligger över den lägsta frekvensen den är konstruerad för (där den
är ungefär en halv våglängd eller kortare) och riktningen följer hur
tråden ligger. En longwire mot sydväst gynnar sydväst, med flera lober
snarare än en enda. Antalet lober blir mer och mer markerat ju högre
frekvens man använder antennen på och ju längre den är i förhållande
till våglängden desto större blir riktverkan i trådens riktning.

Vanligare i dag är den ändmatade halvvågen. Samma längd som en dipol,
men matad i änden där fästet sitter. Impedansen där är hög, några
tusen $\Omega$, så den går inte rakt in i radion utan behöver en unun,
en transformator ofta 49:1 som tar ner impedansen mot
50\,$\Omega$. Det behövs också en motvikt. Ibland räcker koaxens skärm
om en strömbalun sitter en bit ner på kabeln, annars en kort tråd vid
matningspunkten.

Till skillnad från en mittmatad dipol tar den även de jämna
övertonerna. En halvvåg på 40\,m går alltså också på 20, 15 och 10\,m.
Den är för kort för 80 och 160\,m. 30, 17 och 12\,m ligger inte i den
kedjan. Där får man ta en annan tråd, eller en tuner. En tråd som inte
ens är en halvvåg på bandet kräver tuner och en motvikt. Annars blir
radion själv en del av antennen, och det märks som radiofrekvens i
rummet. För portabelt är den ändmatade halvvågen ofta det enklaste.
Tråden väger lite, fästet sitter i ena änden och radion står där man
matar.

Fördelen med den ändmatade halvvågen är att den fungerar på alla
resonanta band och rätt utförd behövs inte heller någon
antennanpassare annat än en unun. Nackdelen är att en unun med
omsättningstalet 49:1 också har en del förluster eftersom den är svår
att linda med lite grövre tråd så den blir bra. Omsättningen är ju
alltså 7 gånger så många varv till den ändmatade tråden som till den
drivande delen från radion.

Här finns dock en variant som är väldigt trevlig och det är en så
kallad {\it end fed random wire} eller EFRW som inte alls är så random
som det låter. Jag ska försöka förklara -- om vi väljer en trådlängd
som är icke-resonant på de band vi vill använda den på och samtidigt i
närheten av säg 40\,\% av längden av en halv våglängd så sjunker
matningsimpedansen ordentligt i änden. Nu räcker det ofta med 9:1 i
stället och då blir unun enklare att linda (3 ggr så många varv på
antennsidan som på radiosidan) och vi klarar oss ofta med radions
inbyggda tuner (eller en enkel extern tuner) på många band.

Nu blir impedansen i änden några hundra~$\Omega$ i stället för flera
tusen som tidigare. Det finns en myt om att antenner som inte är i
resonans inte är effektiva men det är trams, det handlar i princip
bara om antennens längd i förhållande till den våglängd man vill
köra. Matematiskt är den bästa dipolen med högst antennvinst den som
är 1,25 våglängder lång.

Detta är en utmärkt portabelantenn eftersom den består av en unun, en
motvikt och en rulle tråd man kan kasta ut där man vill köra.

\begin{figure}[H]
\centering
\begin{tikzpicture}[
  x=1cm, y=1cm,
  antenn/.style={rhbrav, line width=1.35pt, line cap=round},
  motvikt/.style={rhbtack!85!black, line width=1.1pt, line cap=round},
  koax/.style={black!80, line width=1.1pt, line cap=round},
  lina/.style={black!45, line width=0.6pt, dash pattern=on 2pt off 1.5pt},
  isolator/.style={fill=white, draw=black!75, line width=0.65pt, rounded corners=1pt},
  matt/.style={{Stealth[length=1.6mm]}-{Stealth[length=1.6mm]}, black!65, line width=0.55pt},
  nod/.style={fill=black, circle, inner sep=1.2pt}
]

  % ====================================================================
  % VÄNSTER DEL: EFHW (RESONANT HALVVÅG MED 49:1 UNUN)
  % ====================================================================
  \begin{scope}[shift={(-3.45, 0)}]
    \node[etikett, font=\footnotesize\bfseries] at (0, 3.10) {EFHW (resonant, 49:1)};

    % Mark och jord
    \fill[rhbjord] (-3.15, -2.40) rectangle (3.15, -2.00);
    \draw[marklinje] (-3.15, -2.00) -- (3.15, -2.00);

    % Låg stolpe vid matningsänden + hög stolpe/träd i bortre änden
    \draw[black!70, line width=1.3pt, line cap=round] (-2.45, -2.00) -- (-2.45, 0.45);
    \draw[black!70, line width=1.4pt, line cap=round] ( 2.85, -2.00) -- ( 2.85, 2.05);

    % 49:1 Unun vid matningspunkten
    \filldraw[fill=rhbjord!85, draw=black!75, line width=0.65pt, rounded corners=1.5pt]
      (-2.35, -0.15) rectangle (-1.45, 0.45);
    \node[etikett, font=\scriptsize\bfseries] at (-1.90, 0.22) {Unun};
    \node[etikett, font=\scriptsize] at (-1.90, -0.01) {49:1 ($1{:}7$)};

    % Resonant halvvågstråd snett uppåt (hög impedans i änden, ~2500 ohm)
    \draw[antenn] (-1.45, 0.25) -- (2.15, 1.75);
    \node[nod, fill=rhbrav] at (-1.45, 0.25) {};
    \filldraw[isolator, rotate around={22.6:(2.28,1.80)}] (2.13, 1.71) rectangle (2.43, 1.89);
    \draw[lina] (2.41, 1.86) -- (2.85, 1.95);

    % Spännings- och strömkurva (spänningsmax i ändarna -> 2500 ohm)
    \draw[rhbrav!75, line width=0.75pt, dash pattern=on 2.5pt off 1.5pt]
      (-1.45, 0.25) .. controls (-0.15, 1.65) and (0.85, 2.15) .. (2.15, 1.75);
    \node[etikett, text=rhbrav, rotate=22.6, anchor=south] at (0.20, 1.38)
      {$L = \lambda/2$ (eller $n\cdot\lambda/2$)};
    \node[etikett, text=rhbes, anchor=west] at (-1.35, 0.02) {$Z \approx 2\,500\,\Omega$};

    % Kort motvikt (ca 0,05 lambda) ELLER koaxskärm fram till 1:1-balun
    \draw[motvikt] (-2.15, -0.15) -- (-2.95, -0.85);
    \node[etikett, text=rhbtack!85!black, anchor=east] at (-2.25, -0.68) {Motvikt};

    % Koax 50 ohm + 1:1 mantelströmsbalun en bit ner + Radio (ingen tuner)
    \draw[koax] (-1.75, -0.15) -- (-1.75, -0.95);
    \filldraw[fill=rhbe!45, draw=black!75, line width=0.6pt, rounded corners=1pt]
      (-2.05, -1.32) rectangle (-1.45, -0.95);
    \node[etikett, anchor=west] at (-1.38, -1.14) {1:1-balun};
    \draw[koax] (-1.75, -1.32) -- (-1.75, -1.72) -- (0.35, -1.72);
    \node[etikett, anchor=south] at (-0.70, -1.68) {Koax $50\,\Omega$};

    \filldraw[fill=white, draw=black!75, line width=0.65pt, rounded corners=1.5pt]
      (0.35, -1.95) rectangle (1.65, -1.45);
    \node[etikett, font=\scriptsize\bfseries] at (1.00, -1.70) {Radio};
  \end{scope}

  % ====================================================================
  % HÖGER DEL: EFRW ("RANDOM WIRE", ICKE-RESONANT MED 9:1 UNUN + TUNER)
  % ====================================================================
  \begin{scope}[shift={(3.45, 0)}]
    \node[etikett, font=\footnotesize\bfseries] at (0, 3.10) {EFRW (icke-resonant, 9:1)};

    % Mark och jord
    \fill[rhbjord] (-3.15, -2.40) rectangle (3.15, -2.00);
    \draw[marklinje] (-3.15, -2.00) -- (3.15, -2.00);

    % Låg stolpe vid matningsänden + hög stolpe/träd i bortre änden
    \draw[black!70, line width=1.3pt, line cap=round] (-2.45, -2.00) -- (-2.45, 0.45);
    \draw[black!70, line width=1.4pt, line cap=round] ( 2.85, -2.00) -- ( 2.85, 2.05);

    % 9:1 Unun vid matningspunkten
    \filldraw[fill=rhbjord!85, draw=black!75, line width=0.65pt, rounded corners=1.5pt]
      (-2.35, -0.15) rectangle (-1.45, 0.45);
    \node[etikett, font=\scriptsize\bfseries] at (-1.90, 0.22) {Unun};
    \node[etikett, font=\scriptsize] at (-1.90, -0.01) {9:1 ($1{:}3$)};

    % Icke-resonant tråd snett uppåt (måttlig impedans, några hundra ohm)
    \draw[antenn] (-1.45, 0.25) -- (2.15, 1.75);
    \node[nod, fill=rhbrav] at (-1.45, 0.25) {};
    \filldraw[isolator, rotate around={22.6:(2.28,1.80)}] (2.13, 1.71) rectangle (2.43, 1.89);
    \draw[lina] (2.41, 1.86) -- (2.85, 1.95);

    \node[etikett, text=rhbrav, rotate=22.6, anchor=south] at (0.35, 1.05)
      {Icke-resonant längd ($L \neq n\cdot\lambda/2$)};
    \node[etikett, text=rhbes, anchor=west] at (-1.35, 0.02) {$Z \approx 300\text{--}600\,\Omega$};

    % Motvikt utlagd vid matningspunkten
    \draw[motvikt] (-2.10, -0.15) -- (-2.95, -1.85);
    \node[etikett, text=rhbtack!85!black, anchor=east] at (-2.45, -1.05) {Motvikt};

    % Koax 50 ohm + Antenntuner + Radio
    \draw[koax] (-1.75, -0.15) -- (-1.75, -1.72) -- (-0.45, -1.72);
    \node[etikett, anchor=south] at (-1.10, -1.68) {Koax $50\,\Omega$};

    \filldraw[fill=rhbfsvag!75, draw=black!75, line width=0.65pt, rounded corners=1.5pt]
      (-0.45, -1.95) rectangle (0.65, -1.45);
    \node[etikett, font=\scriptsize\bfseries] at (0.10, -1.70) {Tuner};

    \draw[koax] (0.65, -1.70) -- (1.15, -1.70);

    \filldraw[fill=white, draw=black!75, line width=0.65pt, rounded corners=1.5pt]
      (1.15, -1.95) rectangle (2.35, -1.45);
    \node[etikett, font=\scriptsize\bfseries] at (1.75, -1.70) {Radio};
  \end{scope}

\end{tikzpicture}
\caption{Två filosofier för ändmatade trådantenner. Till vänster: \textbf{EFHW}
(ändmatad halvvåg) är avstämd till $\lambda/2$ och dess övertoner. Ändimpedansen
är mycket hög ($\approx 2\,500\,\Omega$) och transformeras ner med en 49:1-unun
($1{:}7$ i varvtal) så att den kan köras utan tuner på sina resonansband. Till
höger: \textbf{EFRW} (\textit{random wire}) klipps medvetet till en icke-resonant
längd där ändimpedansen är några hundra ohm, vilket klarar sig med en enklare
och förlustsnålare 9:1-unun ($1{:}3$) tillsammans med motvikt och antenntuner.}
\label{fig:efhw-efrw}
\end{figure}

\subsection{Yagi-Uda}
\index{Yagi-Uda}
\index{antenn!Yagi-Uda}
\index{riktantenn}
\index{reflektor}
\index{direktor}
\index{antennvinst}
\index{dBd}
\index{dBi}
\index{riktantenner}
\index{antennvinst!Yagi}

En yagi-uda är en dipol med gott sällskap. Bakom den drivna dipolen
sitter en reflektor, något längre. Framför sitter en eller flera
direktorer, något kortare. Reflektorn trycker energin
framåt. Direktorerna samlar den.  Fler element ger mer antennvinst och
smalare lob. Hur vinsten räknas, i dBd och dBi, står i
teknikkapitlet. Dipolen är där 0\,dBd och 2,15\,dBi. Yagin har oftast
mycket högre antennvinster. Framför allt på VHF och UHF kan det bli 10
dBi eller mer. På ännu högre frekvenser kan det bli mer än så.

På kortvåg blir en yagi dock väldigt stor. Tre element på 20\,m har
element kring tio meter och en bom på flera meter. På 15 och 10\,m
krymper det. På 6\,m och uppåt är yagin det vanliga sättet att få
riktverkan, och där får den plats på en vanlig mast.

\begin{figure}[H]
\centering
\begin{tikzpicture}[
  x=1cm, y=1cm,
  bom/.style={black!60, line width=2.2pt, line cap=round},
  element/.style={rhbrav, line width=1.6pt, line cap=round},
  parasit/.style={black!80, line width=1.45pt, line cap=round},
  koax/.style={black!80, line width=1.1pt, line cap=round},
  isolator/.style={fill=white, draw=black!75, line width=0.65pt, rounded corners=1pt},
  matt/.style={{Stealth[length=1.6mm]}-{Stealth[length=1.6mm]}, black!65, line width=0.55pt},
  pil/.style={-{Stealth[length=2.2mm]}, rhbes, line width=1.2pt}
]

  % ====================================================================
  % VÄNSTER DEL: 3-ELEMENTS YAGI-UDA OVANIFRÅN (KONSTRUKTION)
  % ====================================================================
  \begin{scope}[shift={(-2.65, 0)}]
    \node[etikett, font=\footnotesize\bfseries] at (0.15, 2.95) {3-elements Yagi-Uda ovanifrån};

    % Bom (horisontellt bärrör) och mastfäste
    \draw[bom] (-2.85, 0) -- (2.95, 0);
    \filldraw[fill=black!25, draw=black!75, line width=0.7pt] (0, 0) circle[radius=0.18];
    \node[etikett, anchor=north east, text=black!70] at (-0.12, -0.18) {Bom / mast};

    % 1) Reflektor (bakom, längst: ca 0,50 lambda) vid x = -2.20
    \draw[parasit] (-2.20, -2.15) -- (-2.20, 2.15);
    \node[etikett, font=\scriptsize\bfseries, anchor=south] at (-2.20, 2.20) {Reflektor};
    \node[etikett, anchor=north] at (-2.20, -2.20) {Längre ($\approx 0{,}50\lambda$)};

    % 2) Drivet element / matad dipol (i mitten: ca 0,47 lambda) vid x = -0.35
    \draw[element] (-0.35, -1.90) -- (-0.35, -0.14);
    \draw[element] (-0.35,  0.14) -- (-0.35,  1.90);
    \filldraw[isolator] (-0.45, -0.14) rectangle (-0.25, 0.14);
    \node[etikett, text=rhbrav, font=\scriptsize\bfseries, anchor=south] at (-0.35, 2.20) {Drivet element};
    \node[etikett, text=rhbrav, anchor=north] at (-0.35, -2.20) {Matad dipol ($\approx 0{,}47\lambda$)};

    % Balun + Koax 50 ohm från det drivna elementet
    \filldraw[fill=rhbjord!85, draw=black!75, line width=0.55pt, rounded corners=1pt]
      (-0.92, -0.45) rectangle (-0.48, -0.12);
    \draw[koax] (-0.70, -0.45) .. controls (-0.70, -1.15) and (-0.20, -1.25) .. (0.00, -1.65);
    \node[etikett, anchor=west] at (0.05, -1.55) {Koax $50\,\Omega$};

    % 3) Direktor 1 (framför, kortare: ca 0,44 lambda) vid x = +1.35
    \draw[parasit] (1.35, -1.68) -- (1.35, 1.68);
    \node[etikett, font=\scriptsize\bfseries, anchor=south] at (1.35, 2.20) {Direktor 1};

    % 4) Direktor 2 (längst fram, kortast: ca 0,42 lambda) vid x = +2.65
    \draw[parasit] (2.65, -1.55) -- (2.65, 1.55);
    \node[etikett, font=\scriptsize\bfseries, anchor=south] at (2.65, 2.20) {Direktor 2};
    \node[etikett, anchor=north] at (2.00, -2.20) {Kortare ($\approx 0{,}42\text{--}0{,}45\lambda$)};

    % Pil för strålningsriktning (framåt)
    \draw[pil] (2.95, 0) -- (3.85, 0);
    \node[etikett, text=rhbes, font=\scriptsize\bfseries, anchor=south] at (3.40, 0.08) {Framåt};
  \end{scope}

  % ====================================================================
  % HÖGER DEL: RIKTVERKAN OCH STRÅLNINGSLOB (JÄMFÖRT MED DIPOL)
  % ====================================================================
  \begin{scope}[shift={(4.25, 0)}]
    \node[etikett, font=\footnotesize\bfseries] at (0.35, 2.95) {Strålningsmönster};

    % Hjälpcirkel och axlar
    \draw[black!20, line width=0.5pt] (0,0) circle[radius=1.75];
    \draw[black!30, line width=0.45pt, dash pattern=on 2pt off 2pt] (-1.65, 0) -- (2.65, 0);
    \draw[black!30, line width=0.45pt, dash pattern=on 2pt off 2pt] (0, -2.00) -- (0, 2.00);

    % Vanlig halvvågsdipol som referens (0 dBd = 2,15 dBi, streckad åtta)
    \draw[black!40, line width=0.65pt, dash pattern=on 2pt off 1.5pt]
      (0,0)
      .. controls ( 0.40,  0.85) and ( 1.45,  0.90) .. ( 1.45, 0)
      .. controls ( 1.45, -0.90) and ( 0.40, -0.85) .. (0,0)
      .. controls (-0.40,  0.85) and (-1.45,  0.90) .. (-1.45, 0)
      .. controls (-1.45, -0.90) and (-0.40, -0.85) .. (0,0);
    \node[etikett, text=black!60, anchor=east] at (-0.85, 1.15) {Dipol ($0\text{ dBd}$)};

    % Yagi-Uda: stark huvudlob framåt + små sido- och baklober
    \filldraw[fill=rhbtack!25, draw=rhbtack!85!black, line width=1.1pt]
      (0,0)
      .. controls ( 0.55,  0.95) and ( 2.55,  1.10) .. ( 2.55,  0.00)
      .. controls ( 2.55, -1.10) and ( 0.55, -0.95) .. ( 0.00,  0.00);

    % Liten baklob (fram/bak-förhållande F/B) och sidolober
    \filldraw[fill=rhbtack!18, draw=rhbtack!85!black, line width=0.85pt]
      (0,0)
      .. controls (-0.25,  0.22) and (-0.65,  0.25) .. (-0.65,  0.00)
      .. controls (-0.65, -0.25) and (-0.25, -0.22) .. ( 0.00,  0.00);
    \filldraw[fill=rhbtack!18, draw=rhbtack!85!black, line width=0.75pt]
      (0,0)
      .. controls ( 0.15,  0.45) and (-0.15,  0.45) .. ( 0.00,  0.00)
      .. controls ( 0.15, -0.45) and (-0.15, -0.45) .. ( 0.00,  0.00);

    \node[etikett, text=rhbtack!85!black, font=\scriptsize\bfseries, anchor=south] at (1.55, 0.82)
      {Huvudlob (hög vinst)};
    \node[etikett, text=black!70, anchor=north] at (-0.65, -0.32) {Dämpad baklob};
  \end{scope}

\end{tikzpicture}
\caption{Principen för en Yagi-Uda-antenn. Endast ett element (dipolen) matas
med koaxialkabel. Bakom dipolen sitter en något längre reflektor som speglar
energin framåt, och framför sitter en eller flera något kortare direktorer som
samlar strålningen i en smal huvudlob med hög antennvinst (mätt i dBd eller dBi).}
\label{fig:yagi}
\end{figure}

\subsection{Polarisation}
\index{polarisation}
\index{antenn!polarisation}
\index{horisontell polarisation}
\index{vertikal polarisation}
\index{polarisation!rymdvåg}
\index{polarisation!fri sikt}
\index{repeater!polarisation}

En horisontell tråd ger horisontell polarisation. Ett stående spröt
ger vertikal. På rymdvåg genom jonosfären blandas polarisationen, så
en dipol och en vertikal kan ändå höra varandra. På markvåg och på fri
sikt stannar den. Där ska båda köra samma, annars försvinner en stor
del av signalen, uppåt 20\,dB i värsta fall.

FM, repeatertrafik och det mesta mobila är vertikalt. SSB och CW på
6\,m, 2\,m och 70\,cm är som regel horisontellt. En yagi monteras
efter det man vill köra. Elementen horisontellt för SSB, stående för
repeater.

\subsection{VHF och 70 cm}
\label{sec:vhf-antenn}
\index{antenn!VHF}
\index{antenn!70 cm}
\index{antenn!6 m}
\index{gummiantenn}
\index{handapparat!antenn}
\index{repeater}
\index{ERP}
\index{2 m!antenn}
\index{70 cm}
\index{kvartsvåg!VHF}
\index{kabeldämpning}

På 6\,m är en halvvåg omkring tre meter, på 2\,m omkring en meter och
på 70\,cm en tredjedel av det. En liten yagi är därför det vanliga
sättet att få räckvidd, och en vertikal på taket räcker långt mot
repeatrar. Höjd slår effekt. Samma radio hör en repeater med antennen
över taknocken och är tyst när den sitter i fönstret.

Riktningen spelar roll så fort det är tropo, sporadiskt E eller
norrsken. Fri sikt är snällare. Där räcker det att komma över hindren.

Gummiantennen på handapparaten är en kompromiss. Den är kortad och
förlustrik, och radions chassi får vara jordplan. Håller man handen
över den blir den ännu sämre. En kvartsvåg på bilen, eller en yagi på
masten, är något annat. På 70\,cm straffar kabeln också. Lång tunn
koax mellan radio och antenn kan äta mer effekt än antennen
vinner. Det är samma kabeldämpning som räknas i teknikkapitlet.

På 2\,m, där många kör med instegsbehörighet, räknas gränsen som
ERP. En yagi har vinst över dipolen, så sändareffekten får sänkas.
Hur det räknas står i avsnitt~\ref{sec:insteg-restriktioner}. Vilka
band som är öppna för instegscertifikat står i
avsnitt~\ref{sec:instegsband}.

\subsection{Portabelt}
\label{sec:portabelt}
\index{portabel}
\index{antenn!portabel}
\index{portabel!endmatad}
\index{portabel!inverterad V}
\index{portabel!vertikal}
\index{portabel!VHF}
\index{portabel!yagi}
\index{portabel!kvartsvåg}

Ute räknas vikt, ett fäste och hur fort tråden är i luften. Tre
antenner klarar det bättre än de andra.

Den ändmatade halvvågen är den man oftast tar med. En rulle tråd, en
transformator 49:1 och ett rep i ett träd eller ett spö. Matningen
sitter där radion står. En halvvåg på 40\,m täcker också 20, 15 och
10\,m.

Inverterad V är nästan lika bra när man har ett spö att sätta i mitten.
Ingen transformator utöver en strömbalun 1:1, bara koax. Ändarna får
hänga lägre, så länge ingen når dem.

Vertikalen med tre eller fyra upphöjda radialer är valet på öppen
mark, vid en strand eller där träden saknas. Ett spö, och jordplanet
hänger med. Över salt vatten är den svår att slå.

På VHF och 70\,cm är det en liten yagi, eller en kvartsvåg på samma
spö. Gummiantennen på handapparaten är vad man har i fickan.
Räckvidden blir kortare.

En rak dipol vill ha två höga fästen. Kortvågsyagi, doublet med stege
och off-center-dipol med balunen hängande i spannet är tyngre att få
upp på en kvart. De går att ta med. De är sällan det snabbaste sättet
att komma i luften. I krislådan är det samma avvägning, se
avsnitt~\ref{sec:krislada} i bilaga~\ref{sec:prepping}.

\subsection{Matning och anpassning}
\index{balun}
\index{strömbalun}
\index{antenn!matning}
\index{koaxialkabel!antenn}
\index{antenntuner}
\index{stege}
\index{stående våg}
\index{VSWR}
\index{return loss}
\index{impedans}
\index{koaxialkabel}
\index{anpassning}
\index{missanpassning}
\index{balun!1:1}
\index{strömbalun!1:1}

Radion och den vanliga koaxen är gjorda för 50\,$\Omega$. Dipolen
ligger i närheten, och vertikalen med rätt radialer också. Då kan man
köra utan tuner. En strömbalun 1:1 vid dipolens matningspunkt är ändå
värd att ha. Den håller strömmen på skärmen nere, så att kabeln inte
strålar och loben blir något annat än man räknat med.

En antenntuner vid radion lugnar slutsteget. Effekten som studsar fram
och tillbaka i koaxen är kvar som förlust. På kortvåg, med kort och
grov kabel, är den förlusten ofta måttlig. På VHF växer den fort. Ska
en och samma dipol köras på flera band är den stegmatade doubletten
rätt väg.  Den är beskriven ovan. Tunern ska sitta där stegen kommer
in.

Hur stående våg och return loss hänger ihop, och vad som är en bra
respektive dålig anpassning, står i teknikkapitlet under
missanpassning.






\section{Impedanstransformatorer}
\label{sec:baluner}
\index{balun}
\index{unun}
\index{impedansanpassning}
\index{spänningsbalun}
\index{strömbalun}

En BALUN {\it Balanced-Unbalanced} sitter mellan en symmetrisk antenn
och den osymmetriska koaxen. En unun {\it Unbalanced-Unbalanced}
sitter mellan två osymmetriska sidor med gemensam jord, och den finns
till för att ändra impedansen. När omsättningen inte är 1:1 står
antennens impedans före kolon och de 50\,$\Omega$ som koaxen vill se
efter, samma sätt att skriva som redan används för 4:1, 9:1 och 49:1 i
antennavsnitten.

Det som avgör om matarledningen börjar stråla och bete sig är graden
av common mode-ström som går i den, inte siffran i omsättningen. En
spänningsbalun av Ruthroff-typ tvingar fram spänningar mot jord. En
strömbalun av Guanella-typ släpper i praktiken bara igenom ström som
går åt motsatt håll i de två ledarna, den är ett filter för common
mode-ström, på engelska \textit{common mode current}, där strömmen då
möter ferriten som ett växelströmsmotstånd. Dipolen, den inverterade
V-dipolen och den stegmatade doubletten vill ha strömbalunen. Den
off-center-matade dipolen vill ha en Guanella 4:1, och den ändmatade
tråden en unun.

\begin{table}[H]
\centering
\small
\begin{tabular}{lcll}
  \textbf{Anpassare} & \textbf{Omsättning} & \textbf{Tar} & \textbf{Hör hemma vid} \\ \hline
        Guanella & 1:1  & 50\,$\Omega$ mot 50\,$\Omega$       & Klassisk dipol, inverterad V, doublet (vid tuner) \\
	Ruthroff & 1:1  & 50\,$\Omega$ mot 50\,$\Omega$       & Symmetrisk dipol (spänningsbalun)                 \\
	Ruthroff & 4:1  & 200\,$\Omega$ mot 50\,$\Omega$      & 200\,$\Omega$-matning (helst med 1:1-drossel)     \\
	Ruthroff & 6:1  & 300\,$\Omega$ mot 50\,$\Omega$      & Vikt dipol eller 300\,$\Omega$-ledning            \\
	Guanella & 4:1  & 200\,$\Omega$ mot 50\,$\Omega$      & OCFD, linda på två kärnor, inte en                \\
	UNUN     & 9:1  & 450\,$\Omega$ mot 50\,$\Omega$      & Ändmatad icke-resonant tråd, EFRW                 \\
	UNUN     & 49:1 & ca 2\,500 $\Omega$ mot 50\,$\Omega$ & Ändmatad resonant tråd, EFHW                      \\
\end{tabular}
\caption{Översikt över olika typer och deras användningsområden}
\label{tab:balunval}
\end{table}

\subsection{Balun}
\label{sec:balun}

\subsubsection{Strömbalun 1:1, Guanella}
\label{sec:guanella-11}
\index{strömbalun!1:1}
\index{balun!1:1}
\index{Guanella!1:1}
\index{mantelström}

Strömbalunen 1:1 ändrar inte impedansen. Två ledare, en bifilär
(tvåtrådig) lindning eller en bit koax, ligger på en ferritkärna och
går rakt igenom utan att korsas. Ström som går ut i den ena ledaren
och tillbaka i den andra ser bara transmissionsledningen. Ström som
vill gå åt samma håll i båda, ut på skärmen och tillbaka någon
annanstans, får magnetisera kärnan och dämpas.

Den hör hemma i matningspunkten på en mittmatad dipol och på en
inverterad V, även när impedansen redan ligger nära 50\,$\Omega$. När
stegen från en doublet ska över till en kort koaxstump vid en vanlig
tuner sitter samma 1:1 alldeles intill tunern. En 4:1 är fel verktyg
där. Stegens impedans är inte 200\,$\Omega$, och tunern ska se det
stegen faktiskt lämnar.

 \begin{figure}[H]
    \centering
    \begin{tikzpicture}[
      x=1cm, y=1cm,
      ledareA/.style={rhbrav, line width=1.15pt, line cap=round, line join=round},
      ledareB/.style={rhbes, line width=1.15pt, line cap=round, line join=round},
      ledareAbak/.style={rhbrav!45, line width=0.9pt, line cap=round, dash pattern=on 2pt off 1.5pt},
      ledareBbak/.style={rhbes!45, line width=0.9pt, line cap=round, dash pattern=on 2pt off 1.5pt},
      koppling/.style={black!80, line width=0.8pt, line cap=round, line join=round},
      nod/.style={fill=black, circle, inner sep=1.3pt}
    ]

      % ====================================================================
      % VÄNSTER DEL: PRINCIPSCHEMA (RAKT IGENOM UTAN KORSNING)
      % ====================================================================
      \node[etikett, font=\footnotesize\bfseries] at (-3.6, 2.75) {Principschema};

      % Ferritkärna (horisontella dubbelstreck mellan lindning A och B)
      \draw[black!65, line width=1.1pt] (-3.95, 0.32) -- (-2.35, 0.32);
      \draw[black!65, line width=1.1pt] (-3.95, 0.16) -- (-2.35, 0.16);

      % Övre lindning A (blå, A1 -> A2)
      \draw[ledareA] (-3.95, 0.65)
        arc (180:0:0.20 and 0.24)
        arc (180:0:0.20 and 0.24)
        arc (180:0:0.20 and 0.24)
        arc (180:0:0.20 and 0.24);
      \fill[rhbrav] (-3.85, 0.96) circle[radius=0.055]; % Fasprick A1
      \node[etikett, text=rhbrav, anchor=south east] at (-3.92, 0.70) {$A_1$};
      \node[etikett, text=rhbrav, anchor=south west] at (-2.38, 0.70) {$A_2$};

      % Undre lindning B (rödbrun, B1 -> B2)
      \draw[ledareB] (-3.95, -0.17)
        arc (180:360:0.20 and 0.24)
        arc (180:360:0.20 and 0.24)
        arc (180:360:0.20 and 0.24)
        arc (180:360:0.20 and 0.24);
      \fill[rhbes] (-3.85, -0.48) circle[radius=0.055]; % Fasprick B1
      \node[etikett, text=rhbes, anchor=north east] at (-3.92, -0.22) {$B_1$};
      \node[etikett, text=rhbes, anchor=north west] at (-2.38, -0.22) {$B_2$};

      % Pilar för differentiell ström (lika stor, motsatt riktning)
      \draw[-{Stealth[length=1.6mm]}, rhbrav, line width=0.65pt] (-3.55, 1.15) -- (-2.75, 1.15)
        node[midway, above, etikett, text=rhbrav] {$I$};
      \draw[{Stealth[length=1.6mm]}-, rhbes, line width=0.65pt] (-3.55, -0.67) -- (-2.75, -0.67)
        node[midway, below, etikett, text=rhbes] {$I$};

      % Koaxialkabel 50 ohm (vänster)
      \fill[rhbjord!70, draw=black!70, line width=0.6pt] (-6.15, 0.35) rectangle (-5.05, 0.95);
      \draw[black!70, line width=0.6pt] (-5.05, 0.65) ellipse[x radius=0.08, y radius=0.30];
      \node[etikett, anchor=south] at (-5.60, 1.02) {Koax $50\,\Omega$};

      % Mittledare rakt till A1
      \draw[koppling] (-5.05, 0.65) -- (-3.95, 0.65);
      \node[etikett, anchor=south] at (-4.50, 0.68) {Mittledare};

      % Skärm rakt till B1
      \draw[koppling] (-5.25, 0.35) -- (-5.25, -0.17) -- (-3.95, -0.17);
      \node[nod] at (-5.25, 0.35) {};
      \node[etikett, anchor=north] at (-4.60, -0.22) {Skärm};

      % A2 rakt till Antennben 1
      \draw[koppling] (-2.35, 0.65) -- (-1.15, 0.65);
      \draw[black, fill=white, line width=0.7pt] (-1.15, 0.65) circle[radius=0.07];
      \node[etikett, anchor=west] at (-1.00, 0.65) {Antennben 1};

      % B2 rakt till Antennben 2
      \draw[koppling] (-2.35, -0.17) -- (-1.15, -0.17);
      \draw[black, fill=white, line width=0.7pt] (-1.15, -0.17) circle[radius=0.07];
      \node[etikett, anchor=west] at (-1.00, -0.17) {Antennben 2};

      % Impedansmarkering 50 ohm på antennsidan
      \draw[{Stealth[length=1.6mm]}-{Stealth[length=1.6mm]}, black!60, line width=0.5pt]
        (-1.55, -0.05) -- (-1.55, 0.53)
        node[midway, left=2pt, etikett] {$\approx 50\,\Omega$};


      % ====================================================================
      % HÖGER DEL: PRAKTISK BIFILÄR LINDNING PÅ TOROIDKÄRNA
      % ====================================================================
      \begin{scope}[shift={(3.2, 0.20)}]
        \node[etikett, font=\footnotesize\bfseries] at (0, 2.55) {Bifilär lindning på toroidkärna};

        \def\Rout{1.65}
        \def\Rin{0.90}

        % Baksida av lindningsvarven (streckade under kärnan)
        \foreach \ang/\nextang in {246/216, 216/186, 186/156, 156/126, 126/96} {
          \draw[ledareAbak] ({\ang-3}:\Rin) -- ({\nextang+3}:\Rout);
          \draw[ledareBbak] ({\ang+3}:\Rin) -- ({\nextang+9}:\Rout);
        }

        % Toroidkärna (ring)
        \fill[rhbjord!75, even odd rule, draw=black!65, line width=0.8pt]
          (0,0) circle[radius=\Rout]
          (0,0) circle[radius=\Rin];

        % Framsida av lindningsvarven (två trådar A och B parallellt genom hålet)
        \foreach \ang in {246, 216, 186, 156, 126, 96} {
          % Tråd A (blå)
          \draw[ledareA] ({\ang+3}:\Rout) -- ({\ang-3}:\Rin);
          % Tråd B (rödbrun, direkt bredvid A)
          \draw[ledareB] ({\ang+9}:\Rout) -- ({\ang+3}:\Rin);
        }

        % Startändar nere (från ytterkant: A1 och B1)
        \coordinate (StartA) at (249:\Rout);
        \coordinate (StartB) at (255:\Rout);

        % Slutändar uppe (från innerkant: A2 och B2)
        \coordinate (SlutA) at (93:\Rin);
        \coordinate (SlutB) at (99:\Rin);

        % --- UPPTILL: ANTENNSIDA (A2 och B2 går rakt till var sitt ben) ---
        \draw[ledareA] (SlutA) .. controls (-0.05, 1.45) and (-0.45, 1.85) .. (-1.25, 1.95);
        \draw[black, fill=white, line width=0.7pt] (-1.25, 1.95) circle[radius=0.07];
        \node[etikett, anchor=east] at (-1.35, 1.95) {Antennben 1};
        \node[etikett, text=rhbrav, anchor=west] at (-0.02, 1.28) {$A_2$};

        \draw[ledareB] (SlutB) .. controls (-0.20, 1.45) and ( 0.35, 1.85) .. ( 1.25, 1.95);
        \draw[black, fill=white, line width=0.7pt] (1.25, 1.95) circle[radius=0.07];
        \node[etikett, anchor=west] at (1.35, 1.95) {Antennben 2};
        \node[etikett, text=rhbes, anchor=east] at (-0.22, 1.28) {$B_2$};

        % --- NEDTILL: KOAXSIDA (A1 till mittledare, B1 till skärm) ---
        \coordinate (KoaxMitt)  at (-0.55, -2.05);
        \coordinate (KoaxSkarm) at (-0.10, -2.05);

        \draw[ledareA] (StartA) -- (-0.55, -2.05);
        \draw[ledareB] (StartB) .. controls (-0.35, -1.75) and (-0.10, -1.85) .. (KoaxSkarm);

        \node[etikett, text=rhbrav, anchor=east] at (-0.65, -1.72) {$A_1$};
        \node[etikett, text=rhbes, anchor=west]  at (-0.22, -1.72) {$B_1$};

        % Koaxsymbol nedtill
        \fill[rhbjord!70, draw=black!70, line width=0.6pt] (-0.80, -2.65) rectangle (0.15, -2.05);
        \node[nod] at (KoaxSkarm) {};
        \node[etikett, anchor=east] at (-0.88, -2.35) {Mittledare ($A_1$)};
        \node[etikett, anchor=west] at ( 0.23, -2.35) {Skärm ($B_1$)};
        \node[etikett, anchor=north] at (-0.32, -2.70) {Koax $50\,\Omega$};
      \end{scope}

    \end{tikzpicture}
    \caption{Guanella 1:1 strömbalun. Två ledare ($A$ och $B$) lindas bifilärt
    runt kärnan och går rakt igenom från koaxen till var sitt antennben utan att
    korsas. Differentiell ström ($I$) går fritt genom ledningen, medan
    mantelström (common mode) bromsas av ferritkärnan.}
    \label{fig:guanella-11}
    \end{figure}


\subsubsection{Ruthroff, 1:1}
\label{sec:ruthroff-11}
\index{Ruthroff!1:1}
\index{spänningsbalun!1:1}
\index{balun!1:1}

Ruthroff 1:1 är också 50\,$\Omega$ till 50\,$\Omega$, men den är en
spänningsbalun. Tre lika lindningar på samma kärna kopplas i serie.
Mittledaren går till toppen, skärmen till uttaget mellan den andra och
den tredje lindningen, och antennens två ben tas från det övre uttaget
och från botten. Benen hamnar då på var sin sida om jord med samma
spänning till beloppet, och impedansen följer med oförändrad.

Undertryckningen är det som saknas. Det är spänningen som hålls, inte
strömmen, så en osymmetrisk antenn eller en koax som inte hänger fritt
driver mantelström igenom ändå. En del av effekten kan värma kärnan i
stället för att gå ut i tråden. Den här ersätter inte strömbalunen i
föregående avsnitt.

 \begin{figure}[H]
    \centering
    \begin{tikzpicture}[
      x=1cm, y=1cm,
      ledareA/.style={rhbrav, line width=1.10pt, line cap=round, line join=round},
      ledareB/.style={rhbes, line width=1.10pt, line cap=round, line join=round},
      ledareC/.style={rhbtack, line width=1.10pt, line cap=round, line join=round},
      ledareAbak/.style={rhbrav!45, line width=0.85pt, line cap=round, dash pattern=on 2pt off 1.5pt},
      ledareBbak/.style={rhbes!45, line width=0.85pt, line cap=round, dash pattern=on 2pt off 1.5pt},
      ledareCbak/.style={rhbtack!45, line width=0.85pt, line cap=round, dash pattern=on 2pt off 1.5pt},
      koppling/.style={black!80, line width=0.8pt, line cap=round, line join=round},
      nod/.style={fill=black, circle, inner sep=1.3pt}
    ]

      % ====================================================================
      % VÄNSTER DEL: PRINCIPSCHEMA (TRE SERIEKOPPLADE LINDNINGAR)
      % ====================================================================
      \node[etikett, font=\footnotesize\bfseries] at (-3.6, 3.20) {Principschema};

      % Ferritkärna (två vertikala linjer längs alla tre lindningarna)
      \draw[black!65, line width=1.1pt] (-3.15, -2.15) -- (-3.15, 2.45);
      \draw[black!65, line width=1.1pt] (-2.98, -2.15) -- (-2.98, 2.45);

      % Övre lindning A (blå, A1 -> A2)
      \draw[ledareA] (-3.35, 2.40)
        arc (90:-90:0.22 and 0.16)
        arc (90:-90:0.22 and 0.16)
        arc (90:-90:0.22 and 0.16)
        arc (90:-90:0.22 and 0.16);
      \fill[rhbrav] (-3.50, 2.27) circle[radius=0.05]; % Fasprick A1
      \node[etikett, text=rhbrav, anchor=east] at (-3.60, 2.27) {$A_1$};
      \node[etikett, text=rhbrav, anchor=east] at (-3.45, 1.25) {$A_2$};

      % Mellersta lindning B (rödbrun, B1 -> B2)
      \draw[ledareB] (-3.35, 0.86)
        arc (90:-90:0.22 and 0.16)
        arc (90:-90:0.22 and 0.16)
        arc (90:-90:0.22 and 0.16)
        arc (90:-90:0.22 and 0.16);
      \fill[rhbes] (-3.50, 0.73) circle[radius=0.05]; % Fasprick B1
      \node[etikett, text=rhbes, anchor=east] at (-3.60, 0.73) {$B_1$};
      \node[etikett, text=rhbes, anchor=east] at (-3.45, -0.30) {$B_2$};

      % Undre lindning C (grön, C1 -> C2)
      \draw[ledareC] (-3.35, -0.68)
        arc (90:-90:0.22 and 0.16)
        arc (90:-90:0.22 and 0.16)
        arc (90:-90:0.22 and 0.16)
        arc (90:-90:0.22 and 0.16);
      \fill[rhbtack] (-3.50, -0.81) circle[radius=0.05]; % Fasprick C1
      \node[etikett, text=rhbtack, anchor=east] at (-3.60, -0.81) {$C_1$};
      \node[etikett, text=rhbtack, anchor=east] at (-3.45, -1.85) {$C_2$};

      % Seriekoppling A2-B1 (övre uttag) och B2-C1 (undre uttag)
      \draw[koppling] (-3.35, 1.12) -- (-3.35, 0.86);
      \node[nod] (Uttag1) at (-3.35, 0.99) {};

      \draw[koppling] (-3.35, -0.42) -- (-3.35, -0.68);
      \node[nod] (Uttag2) at (-3.35, -0.55) {};

      % Koaxialkabel 50 ohm (vänster)
      \fill[rhbjord!70, draw=black!70, line width=0.6pt] (-6.05, 0.35) rectangle (-4.95, 0.95);
      \draw[black!70, line width=0.6pt] (-4.95, 0.65) ellipse[x radius=0.08, y radius=0.30];
      \node[etikett, anchor=south] at (-5.50, 1.02) {Koax $50\,\Omega$};

      % Mittledare till toppen (A1)
      \draw[koppling] (-4.95, 0.65) -- (-4.35, 0.65) -- (-4.35, 2.40) -- (-3.35, 2.40);
      \node[etikett, anchor=south] at (-4.35, 2.44) {Mittledare};

      % Skärm till undre uttaget (B2 + C1)
      \draw[koppling] (-5.15, 0.35) -- (-5.15, -0.55) -- (Uttag2);
      \node[nod] at (-5.15, 0.35) {};
      \node[etikett, anchor=north] at (-4.40, -0.60) {Skärm ($0\,\text{V}$)};

      % Antennben 1 från övre uttaget (A2 + B1)
      \draw[koppling] (Uttag1) -- (-1.40, 0.99);
      \draw[black, fill=white, line width=0.7pt] (-1.40, 0.99) circle[radius=0.07];
      \node[etikett, anchor=west] at (-1.25, 0.99) {Antennben 1};

      % Antennben 2 från botten (C2)
      \draw[koppling] (-3.35, -1.96) -- (-1.40, -1.96);
      \draw[black, fill=white, line width=0.7pt] (-1.40, -1.96) circle[radius=0.07];
      \node[etikett, anchor=west] at (-1.25, -1.96) {Antennben 2};

      % Impedansmarkering 50 ohm på antennsidan
      \draw[{Stealth[length=1.6mm]}-{Stealth[length=1.6mm]}, black!60, line width=0.5pt]
        (-1.80, -1.80) -- (-1.80, 0.84)
        node[midway, fill=white, etikett] {$\approx 50\,\Omega$};


      % ====================================================================
      % HÖGER DEL: PRAKTISK TRIFILÄR LINDNING PÅ TOROIDKÄRNA
      % ====================================================================
      \begin{scope}[shift={(3.25, 0.65)}]
        \node[etikett, font=\footnotesize\bfseries] at (0, 2.55) {Trifilär lindning på toroidkärna};

        \def\Rout{1.85}
        \def\Rin{1.02}

        % Baksida av lindningsvarven (streckade under kärnan, 3 trådar per varv)
        \foreach \ang/\nextang in {195/157, 157/119, 119/81, 81/43, 43/5} {
          \draw[ledareAbak] ({\ang-2}:\Rin)  -- ({\nextang+4}:\Rout);
          \draw[ledareBbak] ({\ang-7}:\Rin)  -- ({\nextang-1}:\Rout);
          \draw[ledareCbak] ({\ang-12}:\Rin) -- ({\nextang-6}:\Rout);
        }

        % Toroidkärna (ring)
        \fill[rhbjord!75, even odd rule, draw=black!65, line width=0.8pt]
          (0,0) circle[radius=\Rout]
          (0,0) circle[radius=\Rin];

        % Framsida av lindningsvarven (tre trådar A, B och C parallellt)
        \foreach \ang in {195, 157, 119, 81, 43, 5} {
          \draw[ledareA] ({\ang+4}:\Rout) -- ({\ang-2}:\Rin);
          \draw[ledareB] ({\ang-1}:\Rout) -- ({\ang-7}:\Rin);
          \draw[ledareC] ({\ang-6}:\Rout) -- ({\ang-12}:\Rin);
        }

        % Startändar nere till vänster (från ytterkant: A1, B1, C1)
        \coordinate (StartA) at (199:\Rout);
        \coordinate (StartB) at (194:\Rout);
        \coordinate (StartC) at (189:\Rout);

        % Slutändar nere till höger (från innerkant: A2, B2, C2)
        \coordinate (SlutA) at ( 3:\Rin);
        \coordinate (SlutB) at (-2:\Rin);
        \coordinate (SlutC) at (-7:\Rin);

        % 1) A1 -> Koax mittledare (längst ner till vänster)
        \draw[ledareA] (StartA) -- (-2.35, -0.60) -- (-2.35, -2.25) -- (-0.55, -2.25) -- (-0.55, -2.60);
        \node[etikett, text=rhbrav, anchor=east] at (-2.42, -0.80) {$A_1$};

        % 2) A2 + B1 -> Antennben 1 (övre uttaget)
        \coordinate (NodeAB) at (-0.65, -1.35);
        \draw[ledareB] (StartB) .. controls (-1.55, -0.95) and (-1.05, -1.35) .. (NodeAB);
        \draw[ledareA] (SlutA)  .. controls ( 0.75, -0.45) and ( 0.05, -1.35) .. (NodeAB);
        \node[nod] at (NodeAB) {};
        \draw[koppling] (NodeAB) -- (-3.10, -1.35);
        \draw[black, fill=white, line width=0.7pt] (-3.10, -1.35) circle[radius=0.07];
        \node[etikett, anchor=east] at (-3.20, -1.35) {Antennben 1};
        \node[etikett, text=rhbrav, anchor=south] at (-0.65, -1.30) {$A_2{+}B_1$};

        % 3) B2 + C1 -> Koax skärm (undre uttaget)
        \coordinate (NodeBC) at (0.10, -1.80);
        \draw[ledareC] (StartC) .. controls (-1.35, -1.45) and (-0.55, -1.80) .. (NodeBC);
        \draw[ledareB] (SlutB)  .. controls ( 0.90, -0.85) and ( 0.55, -1.80) .. (NodeBC);
        \node[nod] at (NodeBC) {};
        \draw[koppling] (NodeBC) -- (0.10, -2.60);
        \node[etikett, text=rhbes, anchor=west] at (0.20, -1.78) {$B_2{+}C_1$};

        % 4) C2 -> Antennben 2 (nere till höger)
        \draw[ledareC] (SlutC) .. controls (1.45, -0.45) and (2.05, -0.95) .. (2.35, -1.35) -- (3.00, -1.35);
        \draw[black, fill=white, line width=0.7pt] (3.00, -1.35) circle[radius=0.07];
        \node[etikett, anchor=west] at (3.10, -1.35) {Antennben 2};
        \node[etikett, text=rhbtack, anchor=south west] at (1.70, -0.80) {$C_2$};

        % Koaxsymbol nedtill
        \fill[rhbjord!70, draw=black!70, line width=0.6pt] (-0.80, -3.20) rectangle (0.35, -2.60);
        \node[nod] at (0.10, -2.60) {};
        \node[etikett, anchor=east] at (-0.88, -2.90) {Mittledare ($A_1$)};
        \node[etikett, anchor=west] at ( 0.43, -2.90) {Skärm ($B_2{+}C_1$)};
        \node[etikett, anchor=north] at (-0.22, -3.25) {Koax $50\,\Omega$};
      \end{scope}

    \end{tikzpicture}
    \caption{Ruthroff 1:1 spänningsbalun. Tre lika trådar ($A$, $B$ och $C$)
    lindas trifilärt åt samma håll och seriekopplas ($A_2$ med $B_1$, samt $B_2$
    med $C_1$). Koaxens mittledare läggs på $A_1$ och skärmen på uttaget
    $B_2{+}C_1$. Antennbenen tas från det övre uttaget ($A_2{+}B_1$) och från
    botten ($C_2$).}
    \label{fig:ruthroff-11}
    \end{figure}


\subsubsection{Ruthroff, 4:1}
\label{sec:ruthroff-41}
\index{Ruthroff!4:1}
\index{spänningsbalun!4:1}
\index{balun!4:1}

Två lika lindningar på samma kärna räcker för 4:1. Mittledaren och det
ena antennbenet möts i toppen, skärmen ligger på mittenuttaget och det
andra benet i botten. Spänningen mellan benen blir den dubbla mot den
koaxen lägger på, och impedansen blir den fyrdubbla, så omkring
200\,$\Omega$ kommer ner mot 50\,$\Omega$. Denna balun lindas ofta
bifilärt (med två trådar), se skissen över hur den konstrueras.

Det är den spänningsbalun som förr satt i nästan varje
off-center-matad dipol. Antennen är sned, med olika långa ben, och en
balun som bara håller spänningarna släpper igenom mantelström. Loben
drar snett. Ska en Ruthroff 4:1 ändå användas hör en strömbalun 1:1
hemma på koaxen strax under den.

  \begin{figure}[H]
    \centering
    \begin{tikzpicture}[
      x=1cm, y=1cm,
      ledareA/.style={rhbrav, line width=1.15pt, line cap=round, line join=round},
      ledareB/.style={rhbes, line width=1.15pt, line cap=round, line join=round},
      ledareAbak/.style={rhbrav!45, line width=0.9pt, line cap=round, dash pattern=on 2pt off 1.5pt},
      ledareBbak/.style={rhbes!45, line width=0.9pt, line cap=round, dash pattern=on 2pt off 1.5pt},
      koppling/.style={black!80, line width=0.8pt, line cap=round, line join=round},
      nod/.style={fill=black, circle, inner sep=1.3pt}
    ]

      % ====================================================================
      % VÄNSTER DEL: PRINCIPSCHEMA
      % ====================================================================
      \node[etikett, font=\footnotesize\bfseries] at (-3.6, 3.05) {Principschema};

      % Ferritkärna (två vertikala linjer)
      \draw[black!65, line width=1.1pt] (-3.15, -1.55) -- (-3.15, 2.15);
      \draw[black!65, line width=1.1pt] (-2.98, -1.55) -- (-2.98, 2.15);

      % Övre lindning A (blå, A1 -> A2)
      \draw[ledareA] (-3.35, 2.10)
        arc (90:-90:0.24 and 0.21)
        arc (90:-90:0.24 and 0.21)
        arc (90:-90:0.24 and 0.21)
        arc (90:-90:0.24 and 0.21);
      \fill[rhbrav] (-3.50, 1.95) circle[radius=0.055]; % Fasprick A1
      \node[etikett, text=rhbrav, anchor=east] at (-3.62, 1.95) {$A_1$};
      \node[etikett, text=rhbrav, anchor=east] at (-3.45, 0.55) {$A_2$};

      % Undre lindning B (rödbrun, B1 -> B2)
      \draw[ledareB] (-3.35, 0.18)
        arc (90:-90:0.24 and 0.21)
        arc (90:-90:0.24 and 0.21)
        arc (90:-90:0.24 and 0.21)
        arc (90:-90:0.24 and 0.21);
      \fill[rhbes] (-3.50, 0.03) circle[radius=0.055]; % Fasprick B1
      \node[etikett, text=rhbes, anchor=east] at (-3.62, 0.00) {$B_1$};
      \node[etikett, text=rhbes, anchor=east] at (-3.45, -1.35) {$B_2$};

      % Mittuttag mellan A2 och B1
      \draw[koppling] (-3.35, 0.42) -- (-3.35, 0.18);
      \node[nod] at (-3.35, 0.30) {};

      % Koaxialkabel 50 ohm (vänster)
      \fill[rhbjord!70, draw=black!70, line width=0.6pt] (-5.95, 0.70) rectangle (-4.85, 1.30);
      \draw[black!70, line width=0.6pt] (-4.85, 1.00) ellipse[x radius=0.08, y radius=0.30];
      \node[etikett, anchor=south] at (-5.40, 1.38) {Koax $50\,\Omega$};

      % Mittledare till A1 + Antennben 1
      \draw[koppling] (-4.85, 1.00) -- (-4.30, 1.00) -- (-4.30, 2.10) -- (-1.45, 2.10);
      \node[nod] at (-3.35, 2.10) {};
      \draw[black, fill=white, line width=0.7pt] (-1.45, 2.10) circle[radius=0.07];
      \node[etikett, anchor=west] at (-1.30, 2.10) {Antennben 1};

      % Skärm till mittuttaget (A2 + B1)
      \draw[koppling] (-5.05, 0.70) -- (-5.05, 0.30) -- (-3.35, 0.30);
      \node[nod] at (-5.05, 0.70) {};
      \node[etikett, anchor=north east] at (-4.05, 0.22) {Skärm};

      % B2 till Antennben 2
      \draw[koppling] (-3.35, -1.50) -- (-1.45, -1.50);
      \draw[black, fill=white, line width=0.7pt] (-1.45, -1.50) circle[radius=0.07];
      \node[etikett, anchor=west] at (-1.30, -1.50) {Antennben 2};

      % Impedansmarkering 200 ohm på antennsidan
      \draw[{Stealth[length=1.6mm]}-{Stealth[length=1.6mm]}, black!60, line width=0.5pt]
        (-1.85, -1.35) -- (-1.85, 1.95)
        node[midway, fill=white, etikett] {$\approx 200\,\Omega$};


      % ====================================================================
      % HÖGER DEL: PRAKTISK LINDNING PÅ TOROIDKÄRNA (BIFILÄRT)
      % ====================================================================
      \begin{scope}[shift={(3.2, 0.55)}]
        \node[etikett, font=\footnotesize\bfseries] at (0, 2.50) {Bifilär lindning på toroidkärna};

        \def\Rout{1.85}
        \def\Rin{1.02}

        % Baksida av lindningsvarven (streckade under kärnan)
        \foreach \ang/\nextang in {195/162, 162/129, 129/96, 96/63, 63/30, 30/-3} {
          \draw[ledareAbak] ({\ang-3}:\Rin) -- ({\nextang+3}:\Rout);
          \draw[ledareBbak] ({\ang-9}:\Rin) -- ({\nextang-3}:\Rout);
        }

        % Toroidkärna (ring)
        \fill[rhbjord!75, even odd rule, draw=black!65, line width=0.8pt]
          (0,0) circle[radius=\Rout]
          (0,0) circle[radius=\Rin];

        % Framsida av lindningsvarven (två trådar A och B parallellt genom hålet)
        \foreach \ang in {195, 162, 129, 96, 63, 30, -3} {
          % Tråd A (blå)
          \draw[ledareA] ({\ang+3}:\Rout) -- ({\ang-3}:\Rin);
          % Tråd B (rödbrun, direkt bredvid A)
          \draw[ledareB] ({\ang-3}:\Rout) -- ({\ang-9}:\Rin);
        }

        % Uttag i startänden (nere till vänster: A1 och B1 från ytterkant)
        \coordinate (StartA) at (198:\Rout);
        \coordinate (StartB) at (192:\Rout);

        % Uttag i slutänden (nere till höger: A2 och B2 från innerkant)
        \coordinate (SlutA) at (-6:\Rin);
        \coordinate (SlutB) at (-12:\Rin);

        % Ledare A1 -> Koax mittledare + Antennben 1 (vänster ut)
        \draw[ledareA] (StartA) -- (-2.35, -0.60) -- (-2.35, -2.05);
        \node[nod] at (-2.35, -1.45) {};
        \draw[koppling] (-2.35, -1.45) -- (-3.25, -1.45);
        \draw[black, fill=white, line width=0.7pt] (-3.25, -1.45) circle[radius=0.07];
        \node[etikett, anchor=east] at (-3.35, -1.45) {Antennben 1};
        \node[etikett, text=rhbrav, anchor=east] at (-2.42, -0.82) {$A_1$};

        % Ihopkoppling av A2 och B1 i mitten (mittuttag -> Koax skärm)
        \coordinate (MittUttag) at (0, -1.65);
        \draw[ledareB] (StartB) .. controls (-1.45, -1.25) and (-0.55, -1.65) .. (MittUttag);
        \draw[ledareA] (SlutA)  .. controls ( 0.85, -0.85) and ( 0.45, -1.65) .. (MittUttag);
        \node[nod] at (MittUttag) {};
        \node[etikett, text=rhbes, anchor=north] at (-1.15, -1.32) {$B_1$};
        \node[etikett, text=rhbrav, anchor=north] at ( 0.72, -1.12) {$A_2$};

        % Ledare B2 -> Antennben 2 (höger ut)
        \draw[ledareB] (SlutB) .. controls (1.35, -0.55) and (2.05, -0.95) .. (2.35, -1.45) -- (3.05, -1.45);
        \draw[black, fill=white, line width=0.7pt] (3.05, -1.45) circle[radius=0.07];
        \node[etikett, anchor=west] at (3.15, -1.45) {Antennben 2};
        \node[etikett, text=rhbes, anchor=south west] at (1.65, -0.85) {$B_2$};

        % Anslutning från koaxialkabel nedtill
        \draw[koppling] (-2.35, -2.05) -- (-0.55, -2.05) -- (-0.55, -2.45);
        \draw[koppling] (MittUttag) -- (0, -2.45);

        % Koaxsymbol nedtill
        \fill[rhbjord!70, draw=black!70, line width=0.6pt] (-0.80, -3.05) rectangle (0.25, -2.45);
        \node[nod] at (0, -2.45) {};
        \node[etikett, anchor=east] at (-0.65, -2.25) {Mittledare};
        \node[etikett, anchor=west] at ( 0.12, -2.25) {Skärm ($A_2{+}B_1$)};
        \node[etikett, anchor=north] at (-0.28, -3.10) {Koax $50\,\Omega$};
      \end{scope}

    \end{tikzpicture}
    \caption{Ruthroff 4:1 spänningsbalun. Två lika trådar ($A$ och $B$) lindas
    bifilärt åt samma håll runt kärnan. Slutet på den första lindningen ($A_2$)
    kopplas ihop med början på den andra ($B_1$) till koaxens skärm. Koaxens
    mittledare och det ena antennbenet läggs på $A_1$, och det andra antennbenet
    på $B_2$.}
    \label{fig:ruthroff-41}
    \end{figure}


\subsubsection{Ruthroff, 6:1}
\label{sec:ruthroff-61}
\index{Ruthroff!6:1}
\index{spänningsbalun!6:1}
\index{balun!6:1}

6:1 är en Ruthroff den också, men koaxen ligger inte över halva lindningen. Hela
lindningen är omkring 2,45 gånger den del som sitter mellan mittledare
och skärm. Roten ur sex är det talet, och impedansen följer kvadraten,
så omkring 300\,$\Omega$ tas mot 50\,$\Omega$. En vikt dipol ligger i
den trakten, och en bandkabel på 300\,$\Omega$ också.

Den off-center-matade dipolen hör inte hit. Den ligger kring 200\,$\Omega$
och vill ha 4:1. Doubletten ska inte ha en 6:1 heller, för där är det
stegen och en strömbalun 1:1 vid tunern som gäller. Jordreferensen sitter
inte mitt emellan antennbenen, så spänningarna mot jord blir inte lika
stora, och mantelströmmen är kvar. En strömbalun 1:1 får ta den om
kopplingen ändå används. Ofta räcker en Guanella 4:1 även mot 300\,$\Omega$,
eftersom 300 mot 200 bara är en stående våg på omkring 1,5.

 \begin{figure}[H]
    \centering
    \begin{tikzpicture}[
      x=1cm, y=1cm,
      ledareA/.style={rhbrav, line width=1.15pt, line cap=round, line join=round},
      ledareB/.style={rhbes, line width=1.15pt, line cap=round, line join=round},
      ledareAbak/.style={rhbrav!45, line width=0.9pt, line cap=round, dash pattern=on 2pt off 1.5pt},
      ledareBbak/.style={rhbes!45, line width=0.9pt, line cap=round, dash pattern=on 2pt off 1.5pt},
      koppling/.style={black!80, line width=0.8pt, line cap=round, line join=round},
      nod/.style={fill=black, circle, inner sep=1.3pt}
    ]

      % ====================================================================
      % VÄNSTER DEL: PRINCIPSCHEMA (OSYMMETRISKT UTTAG 1 : 2,45)
      % ====================================================================
      \node[etikett, font=\footnotesize\bfseries] at (-3.6, 3.15) {Principschema};

      % Ferritkärna (två vertikala linjer)
      \draw[black!65, line width=1.1pt] (-3.15, -1.95) -- (-3.15, 2.30);
      \draw[black!65, line width=1.1pt] (-2.98, -1.95) -- (-2.98, 2.30);

      % Övre lindningsdel A (N1 varv, blå, A1 -> uttag)
      \draw[ledareA] (-3.35, 2.25)
        arc (90:-90:0.24 and 0.19)
        arc (90:-90:0.24 and 0.19)
        arc (90:-90:0.24 and 0.19);
      \fill[rhbrav] (-3.50, 2.12) circle[radius=0.055]; % Fasprick A1
      \node[etikett, text=rhbrav, anchor=east] at (-3.62, 2.10) {$A_1$};
      \node[etikett, text=rhbrav, anchor=east] at (-3.62, 1.45) {$N_1$};

      % Undre lindningsdel B (N2 ≈ 1,45*N1 varv, rödbrun, uttag -> B2)
      \draw[ledareB] (-3.35, 0.85)
        arc (90:-90:0.24 and 0.22)
        arc (90:-90:0.24 and 0.22)
        arc (90:-90:0.24 and 0.22)
        arc (90:-90:0.24 and 0.22)
        arc (90:-90:0.24 and 0.22);
      \fill[rhbes] (-3.50, 0.70) circle[radius=0.055]; % Fasprick B1
      \node[etikett, text=rhbes, anchor=east] at (-3.62, -0.25) {$N_2 \approx 1{,}45\,N_1$};
      \node[etikett, text=rhbes, anchor=east] at (-3.45, -1.70) {$B_2$};

      % Uttag mellan del A och del B (ej i mitten!)
      \draw[koppling] (-3.35, 1.11) -- (-3.35, 0.85);
      \node[nod] (Uttag) at (-3.35, 0.98) {};

      % Koaxialkabel 50 ohm (vänster)
      \fill[rhbjord!70, draw=black!70, line width=0.6pt] (-6.00, 1.25) rectangle (-4.90, 1.85);
      \draw[black!70, line width=0.6pt] (-4.90, 1.55) ellipse[x radius=0.08, y radius=0.30];
      \node[etikett, anchor=south] at (-5.45, 1.92) {Koax $50\,\Omega$};

      % Mittledare till A1 + Antennben 1
      \draw[koppling] (-4.90, 1.55) -- (-4.30, 1.55) -- (-4.30, 2.25) -- (-1.40, 2.25);
      \node[nod] at (-3.35, 2.25) {};
      \draw[black, fill=white, line width=0.7pt] (-1.40, 2.25) circle[radius=0.07];
      \node[etikett, anchor=west] at (-1.25, 2.25) {Antennben 1};

      % Skärm till uttaget (efter N1 varv)
      \draw[koppling] (-5.10, 1.25) -- (-5.10, 0.98) -- (Uttag);
      \node[nod] at (-5.10, 1.25) {};
      \node[etikett, anchor=north east] at (-4.00, 0.90) {Skärm (uttag)};

      % B2 till Antennben 2
      \draw[koppling] (-3.35, -1.85) -- (-1.40, -1.85);
      \draw[black, fill=white, line width=0.7pt] (-1.40, -1.85) circle[radius=0.07];
      \node[etikett, anchor=west] at (-1.25, -1.85) {Antennben 2};

      % Impedansmarkering 300 ohm på antennsidan
      \draw[{Stealth[length=1.6mm]}-{Stealth[length=1.6mm]}, black!60, line width=0.5pt]
        (-1.80, -1.70) -- (-1.80, 2.10)
        node[midway, fill=white, etikett] {$\approx 300\,\Omega$};


      % ====================================================================
      % HÖGER DEL: LINDNING PÅ TOROIDKÄRNA MED UTTAG (T.EX. 5 + 7 = 12 VARV)
      % ====================================================================
      \begin{scope}[shift={(3.25, 0.55)}]
        \node[etikett, font=\footnotesize\bfseries] at (0, 2.60) {Lindning med uttag ($5 + 7 = 12$ varv)};

        \def\Rout{1.85}
        \def\Rin{1.02}

        % Baksida av del A (5 varv, blå, streckade under kärnan)
        \foreach \ang/\nextang in {210/190, 190/170, 170/150, 150/130, 130/110} {
          \draw[ledareAbak] (\ang:\Rin) -- (\nextang:\Rout);
        }

        % Baksida av del B (7 varv, rödbrun, fortsätter åt samma håll runt kärnan)
        \foreach \ang/\nextang in {90/70, 70/50, 50/30, 30/10, 10/-10, -10/-30} {
          \draw[ledareBbak] (\ang:\Rin) -- (\nextang:\Rout);
        }

        % Toroidkärna (ring)
        \fill[rhbjord!75, even odd rule, draw=black!65, line width=0.8pt]
          (0,0) circle[radius=\Rout]
          (0,0) circle[radius=\Rin];

        % Framsida av del A (5 varv på vänstra sidan, blå)
        \foreach \ang in {210, 190, 170, 150, 130} {
          \draw[ledareA] (\ang:\Rout) -- (\ang:\Rin);
        }

        % Framsida av del B (7 varv över toppen och högra sidan, rödbrun)
        \foreach \ang in {90, 70, 50, 30, 10, -10, -30} {
          \draw[ledareB] (\ang:\Rout) -- (\ang:\Rin);
        }

        % Start på del A (nere till vänster: A1)
        \coordinate (StartA) at (210:\Rout);
        % Slut på del A / Start på del B (uttaget uppe vid kl. 11)
        \coordinate (SlutA)  at (110:\Rout);
        \coordinate (StartB) at (90:\Rout);
        \coordinate (TapTop) at (100:2.25);

        % Slut på del B (nere till höger: B2)
        \coordinate (SlutB)  at (-30:\Rin);

        % 1) A1 -> Antennben 1 + Koax mittledare
        \draw[ledareA] (StartA) -- (-2.25, -0.92) -- (-2.25, -2.05);
        \node[nod] at (-2.25, -1.45) {};
        \draw[koppling] (-2.25, -1.45) -- (-3.10, -1.45);
        \draw[black, fill=white, line width=0.7pt] (-3.10, -1.45) circle[radius=0.07];
        \node[etikett, anchor=east] at (-3.20, -1.45) {Antennben 1};
        \node[etikett, text=rhbrav, anchor=east] at (-1.45, -0.25) {$N_1 = 5$ varv};
        \node[etikett, text=rhbrav, anchor=east] at (-2.32, -1.05) {$A_1$};

        % 2) Uttaget efter 5 varv (A2 + B1) -> Koax skärm
        \draw[ledareA] (SlutA) -- (TapTop);
        \draw[ledareB] (TapTop) -- (StartB);
        \node[nod] at (TapTop) {};
        \node[etikett, anchor=south] at (TapTop) {Uttag ($A_2{+}B_1$)};
        % Ledning från uttaget ned genom mitten till koaxens skärm
        \draw[koppling] (TapTop) .. controls (-0.15, 1.35) and (0.0, 0.25) .. (0.0, -2.45);

        % 3) B2 (efter 5 + 7 = 12 varv) -> Antennben 2
        \draw[ledareB] (SlutB) .. controls (1.35, -0.65) and (1.95, -1.05) .. (2.25, -1.45) -- (2.95, -1.45);
        \draw[black, fill=white, line width=0.7pt] (2.95, -1.45) circle[radius=0.07];
        \node[etikett, anchor=west] at (3.05, -1.45) {Antennben 2};
        \node[etikett, text=rhbes, anchor=west] at (0.35, 0.25) {$N_2 = 7$ varv};
        \node[etikett, text=rhbes, anchor=south west] at (1.55, -0.95) {$B_2$};

        % 4) Koaxialkabel nedtill
        \draw[koppling] (-2.25, -2.05) -- (-0.55, -2.05) -- (-0.55, -2.45);
        \fill[rhbjord!70, draw=black!70, line width=0.6pt] (-0.80, -3.05) rectangle (0.25, -2.45);
        \node[nod] at (0.0, -2.45) {};
        \node[etikett, anchor=east] at (-0.88, -2.75) {Mittledare ($A_1$)};
        \node[etikett, anchor=west] at ( 0.33, -2.75) {Skärm (uttag)};
        \node[etikett, anchor=north] at (-0.28, -3.10) {Koax $50\,\Omega$};
      \end{scope}

    \end{tikzpicture}
    \caption{Ruthroff 6:1 balun. Hela lindningen har omkring $\sqrt{6} \approx 2{,}45$
    gånger så många varv som delen $N_1$ mellan mittledare och skärm (här visat med
    $N_1 = 5$ varv före uttaget och $N_2 = 7$ varv efter, totalt $12$ varv där
    $(12/5)^2 = 5{,}76 \approx 6$).}
    \label{fig:ruthroff-61}
    \end{figure}


\subsubsection{Guanella, 4:1}
\label{sec:guanella-41}
\index{Guanella!4:1}
\index{strömbalun!4:1}
\index{balun!4:1}
\index{balun!Guanella}

Guanella 4:1 är två strömbaluner 1:1, var och en på sin kärna, med
ledningar kring 100\,$\Omega$. På koaxsidan ligger de parallellt, och
100\,$\Omega$ parallellt med 100\,$\Omega$ är 50\,$\Omega$. På
antennsidan ligger de i serie, och 100\,$\Omega$ i serie med 100\,$\Omega$
är 200\,$\Omega$. Strömmen i varje ledning tvingas gå åt motsatt håll,
så mantelströmmen möter samma drosselverkan som i 1:1.

De två kärnorna är inte en smaksak. Delar lindningarna en och samma
kärna låser det gemensamma flödet balansen, och undertryckningen som var
skälet att välja en strömbalun försvinner. På den off-center-matade
dipolen är det här den 4:1 som ska upp i matningspunkten. En extra
strömbalun 1:1 på koaxen strax under gör nytta när benen är så olika att
en ensam 4:1 inte håller skärmen tyst.

  \begin{figure}[H]
    \centering
    \begin{tikzpicture}[
      x=1cm, y=1cm,
      ledareA/.style={rhbrav, line width=1.15pt, line cap=round, line join=round},
      ledareB/.style={rhbes, line width=1.15pt, line cap=round, line join=round},
      ledareAbak/.style={rhbrav!45, line width=0.9pt, line cap=round, dash pattern=on 2pt off 1.5pt},
      ledareBbak/.style={rhbes!45, line width=0.9pt, line cap=round, dash pattern=on 2pt off 1.5pt},
      koppling/.style={black!80, line width=0.8pt, line cap=round, line join=round},
      nod/.style={fill=black, circle, inner sep=1.3pt}
    ]

      % ====================================================================
      % VÄNSTER DEL: PRINCIPSCHEMA (PARALLELLT IN, SERIE UT)
      % ====================================================================
      \node[etikett, font=\footnotesize\bfseries] at (-3.7, 3.05) {Principschema};

      % --- Övre 1:1-enhet (Kärna 1) ---
      % Ferritkärna 1 (horisontell dubbelstreck mellan A och B)
      \draw[black!65, line width=1.0pt] (-3.95, 1.38) -- (-2.35, 1.38);
      \draw[black!65, line width=1.0pt] (-3.95, 1.24) -- (-2.35, 1.24);
      \node[etikett, anchor=south] at (-3.15, 2.02) {Kärna 1};

      % Lindning A (övre på kärna 1, blå)
      \draw[ledareA] (-3.95, 1.65)
        arc (180:0:0.20 and 0.22)
        arc (180:0:0.20 and 0.22)
        arc (180:0:0.20 and 0.22)
        arc (180:0:0.20 and 0.22);
      \fill[rhbrav] (-3.85, 1.93) circle[radius=0.05];
      \node[etikett, text=rhbrav, anchor=south east] at (-3.92, 1.68) {$A_1$};
      \node[etikett, text=rhbrav, anchor=south west] at (-2.38, 1.68) {$A_2$};

      % Lindning B (undre på kärna 1, rödbrun)
      \draw[ledareB] (-3.95, 0.97)
        arc (180:360:0.20 and 0.22)
        arc (180:360:0.20 and 0.22)
        arc (180:360:0.20 and 0.22)
        arc (180:360:0.20 and 0.22);
      \fill[rhbes] (-3.85, 0.69) circle[radius=0.05];
      \node[etikett, text=rhbes, anchor=north east] at (-3.92, 0.94) {$B_1$};
      \node[etikett, text=rhbes, anchor=north west] at (-2.38, 0.94) {$B_2$};

      % --- Undre 1:1-enhet (Kärna 2) ---
      % Ferritkärna 2 (separat kärna!)
      \draw[black!65, line width=1.0pt] (-3.95, -0.72) -- (-2.35, -0.72);
      \draw[black!65, line width=1.0pt] (-3.95, -0.86) -- (-2.35, -0.86);
      \node[etikett, anchor=north] at (-3.15, -1.52) {Kärna 2};

      % Lindning C (övre på kärna 2, blå)
      \draw[ledareA] (-3.95, -0.45)
        arc (180:0:0.20 and 0.22)
        arc (180:0:0.20 and 0.22)
        arc (180:0:0.20 and 0.22)
        arc (180:0:0.20 and 0.22);
      \fill[rhbrav] (-3.85, -0.17) circle[radius=0.05];
      \node[etikett, text=rhbrav, anchor=south east] at (-3.92, -0.42) {$C_1$};
      \node[etikett, text=rhbrav, anchor=south west] at (-2.38, -0.42) {$C_2$};

      % Lindning D (undre på kärna 2, rödbrun)
      \draw[ledareB] (-3.95, -1.13)
        arc (180:360:0.20 and 0.22)
        arc (180:360:0.20 and 0.22)
        arc (180:360:0.20 and 0.22)
        arc (180:360:0.20 and 0.22);
      \fill[rhbes] (-3.85, -1.41) circle[radius=0.05];
      \node[etikett, text=rhbes, anchor=north east] at (-3.92, -1.16) {$D_1$};
      \node[etikett, text=rhbes, anchor=north west] at (-2.38, -1.16) {$D_2$};

      % --- Koaxialkabel 50 ohm (parallellkoppling på ingången) ---
      \fill[rhbjord!70, draw=black!70, line width=0.6pt] (-6.35, -0.05) rectangle (-5.35, 0.55);
      \draw[black!70, line width=0.6pt] (-5.35, 0.25) ellipse[x radius=0.08, y radius=0.30];
      \node[etikett, anchor=south] at (-5.85, 0.62) {Koax $50\,\Omega$};

      % Mittledare till A1 och C1 (parallellt)
      \draw[koppling] (-5.35, 0.25) -- (-4.75, 0.25);
      \draw[koppling] (-4.75, 1.65) -- (-4.75, -0.45);
      \draw[koppling] (-4.75, 1.65) -- (-3.95, 1.65);
      \draw[koppling] (-4.75, -0.45) -- (-3.95, -0.45);
      \node[nod] at (-4.75, 0.25) {};

      % Skärm till B1 och D1 (parallellt)
      \draw[koppling] (-5.55, -0.05) -- (-5.55, -0.60) -- (-4.40, -0.60);
      \draw[koppling] (-4.40, 0.97) -- (-4.40, -1.13);
      \draw[koppling] (-4.40, 0.97) -- (-3.95, 0.97);
      \draw[koppling] (-4.40, -1.13) -- (-3.95, -1.13);
      \node[nod] at (-5.55, -0.05) {};
      \node[nod] at (-4.40, -0.60) {};
      \node[etikett, anchor=north] at (-5.05, -0.65) {Skärm};

      % --- Antennsida 200 ohm (seriekoppling på utgången) ---
      % A2 -> Antennben 1
      \draw[koppling] (-2.35, 1.65) -- (-1.15, 1.65);
      \draw[black, fill=white, line width=0.7pt] (-1.15, 1.65) circle[radius=0.07];
      \node[etikett, anchor=west] at (-1.00, 1.65) {Antennben 1};

      % B2 seriekopplas med C2
      \draw[koppling] (-2.35, 0.97) -- (-2.00, 0.97) -- (-2.00, -0.45) -- (-2.35, -0.45);
      \node[nod] at (-2.00, 0.26) {};

      % D2 -> Antennben 2
      \draw[koppling] (-2.35, -1.13) -- (-1.15, -1.13);
      \draw[black, fill=white, line width=0.7pt] (-1.15, -1.13) circle[radius=0.07];
      \node[etikett, anchor=west] at (-1.00, -1.13) {Antennben 2};

      % Impedansmarkering 200 ohm
      \draw[{Stealth[length=1.6mm]}-{Stealth[length=1.6mm]}, black!60, line width=0.5pt]
        (-1.45, -0.98) -- (-1.45, 1.50)
        node[midway, fill=white, etikett] {$\approx 200\,\Omega$};


      % ====================================================================
      % HÖGER DEL: PRAKTISK LINDNING PÅ TVÅ SEPARATA TOROIDKÄRNOR
      % ====================================================================
      \begin{scope}[shift={(3.55, 0.25)}]
        \node[etikett, font=\footnotesize\bfseries] at (0, 2.80) {Bifilär lindning på två kärnor};

        \def\Rout{1.32}
        \def\Rin{0.72}
        \def\cxL{-1.55}
        \def\cxR{ 1.55}

        % --- KÄRNA 1 (VÄNSTER) ---
        \begin{scope}[shift={(\cxL, 0.20)}]
          % Baksida av varven (streckade)
          \foreach \ang/\nextang in {245/212, 212/179, 179/146, 146/113} {
            \draw[ledareAbak] ({\ang-3}:\Rin) -- ({\nextang+3}:\Rout);
            \draw[ledareBbak] ({\ang+3}:\Rin) -- ({\nextang+9}:\Rout);
          }
          % Toroidkärna 1
          \fill[rhbjord!75, even odd rule, draw=black!65, line width=0.8pt]
            (0,0) circle[radius=\Rout]
            (0,0) circle[radius=\Rin];
          \node[etikett] at (0.15, 0) {Kärna 1};
          % Framsida av varven
          \foreach \ang in {245, 212, 179, 146, 113} {
            \draw[ledareA] ({\ang+3}:\Rout) -- ({\ang-3}:\Rin);
            \draw[ledareB] ({\ang+9}:\Rout) -- ({\ang+3}:\Rin);
          }
          \coordinate (A1) at (248:\Rout);
          \coordinate (B1) at (254:\Rout);
          \coordinate (A2) at (110:\Rin);
          \coordinate (B2) at (116:\Rin);
        \end{scope}

        % --- KÄRNA 2 (HÖGER) ---
        \begin{scope}[shift={(\cxR, 0.20)}]
          % Baksida av varven (streckade, spegelvänt mot mitten)
          \foreach \ang/\nextang in {-65/-32, -32/1, 1/34, 34/67} {
            \draw[ledareAbak] ({\ang-3}:\Rin) -- ({\nextang+3}:\Rout);
            \draw[ledareBbak] ({\ang+3}:\Rin) -- ({\nextang+9}:\Rout);
          }
          % Toroidkärna 2
          \fill[rhbjord!75, even odd rule, draw=black!65, line width=0.8pt]
            (0,0) circle[radius=\Rout]
            (0,0) circle[radius=\Rin];
          \node[etikett] at (-0.15, 0) {Kärna 2};
          % Framsida av varven
          \foreach \ang in {-65, -32, 1, 34, 67} {
            \draw[ledareA] ({\ang+3}:\Rout) -- ({\ang-3}:\Rin);
            \draw[ledareB] ({\ang+9}:\Rout) -- ({\ang+3}:\Rin);
          }
          \coordinate (C1) at (-68:\Rout);
          \coordinate (D1) at (-62:\Rout);
          \coordinate (C2) at (64:\Rin);
          \coordinate (D2) at (70:\Rin);
        \end{scope}

        % --- SERIEKOPPLING UPPTILL (ANTENNSIDA 200 OHM) ---
        % A2 -> Antennben 1 (uppe till vänster)
        \draw[ledareA] (A2) -- ++(-0.25, 0.95) -- (-2.55, 1.83);
        \draw[black, fill=white, line width=0.7pt] (-2.55, 1.83) circle[radius=0.07];
        \node[etikett, anchor=south] at (-2.55, 1.93) {Antennben 1};
        \node[etikett, text=rhbrav, anchor=east] at (-1.92, 1.45) {$A_2$};

        % B2 ihop med C2 i mitten upptill (seriekoppling)
        \coordinate (SerieTopp) at (0, 1.83);
        \draw[ledareB] (B2) .. controls (-1.45, 1.55) and (-0.65, 1.83) .. (SerieTopp);
        \draw[ledareA] (C2) .. controls ( 1.45, 1.55) and ( 0.65, 1.83) .. (SerieTopp);
        \node[nod] at (SerieTopp) {};
        \node[etikett, anchor=south] at (0, 1.92) {Serie ($B_2{+}C_2$)};

        % D2 -> Antennben 2 (uppe till höger)
        \draw[ledareB] (D2) -- ++(0.25, 0.95) -- (2.55, 1.83);
        \draw[black, fill=white, line width=0.7pt] (2.55, 1.83) circle[radius=0.07];
        \node[etikett, anchor=south] at (2.55, 1.93) {Antennben 2};
        \node[etikett, text=rhbes, anchor=west] at (1.92, 1.45) {$D_2$};

        % --- PARALLELLKOPPLING NEDTILL (KOAXSIDA 50 OHM) ---
        \coordinate (MittPunkt) at (-0.35, -1.55);
        \coordinate (SkarmPunkt) at ( 0.35, -1.90);

        % A1 och C1 (blå) till Koax mittledare
        \draw[ledareA] (A1) .. controls (-1.85, -1.45) and (-0.95, -1.55) .. (MittPunkt);
        \draw[ledareA] (C1) .. controls ( 1.25, -1.45) and ( 0.25, -1.55) .. (MittPunkt);
        \node[nod] at (MittPunkt) {};
        \node[etikett, text=rhbrav, anchor=east] at (-1.98, -1.25) {$A_1$};
        \node[etikett, text=rhbrav, anchor=west] at ( 1.98, -1.25) {$C_1$};

        % B1 och D1 (rödbrun) till Koax skärm
        \draw[ledareB] (B1) .. controls (-1.55, -1.80) and (-0.45, -1.90) .. (SkarmPunkt);
        \draw[ledareB] (D1) .. controls ( 1.65, -1.80) and ( 0.85, -1.90) .. (SkarmPunkt);
        \node[nod] at (SkarmPunkt) {};
        \node[etikett, text=rhbes, anchor=north east] at (-1.35, -1.42) {$B_1$};
        \node[etikett, text=rhbes, anchor=north west] at ( 1.35, -1.42) {$D_1$};

        % Koaxialkabel nedtill
        \draw[koppling] (MittPunkt) -- (-0.35, -2.35);
        \draw[koppling] (SkarmPunkt) -- ( 0.35, -2.35);
        \fill[rhbjord!70, draw=black!70, line width=0.6pt] (-0.60, -2.90) rectangle (0.60, -2.35);
        \node[nod] at (0.35, -2.35) {};
        \node[etikett, anchor=east] at (-0.45, -2.12) {Mittledare ($A_1{+}C_1$)};
        \node[etikett, anchor=west] at ( 0.45, -2.12) {Skärm ($B_1{+}D_1$)};
        \node[etikett, anchor=north] at (0, -2.95) {Koax $50\,\Omega$};
      \end{scope}

    \end{tikzpicture}
    \caption{Guanella 4:1 strömbalun på två separata kärnor. Varje kärna bär en
    bifilär lindning kring $100\,\Omega$. På koaxsidan nedtill kopplas båda
    lindningarna parallellt ($A_1$ med $C_1$ till mittledaren, $B_1$ med $D_1$
    till skärmen) vilket ger $50\,\Omega$. På antennsidan upptill kopplas de i
    serie ($B_2$ ihop med $C_2$) vilket ger $200\,\Omega$ mellan $A_2$ och $D_2$.}
    \label{fig:guanella-41}
    \end{figure}


\subsection{Unun}
\label{sec:unun}
\index{unun}
\index{transformator}

Unun är engelska för \textit{unbalanced to unbalanced}. Båda sidorna har
jord gemensam, så det finns ingen symmetrisk port att hålla, och
omsättningen är hela vitsen. Mantelströmmen stoppas inte här. Den får en
strömbalun 1:1 längre ner på koaxen ta, på samma sätt som redan är sagt
för den ändmatade halvvågen.

\subsubsection{Unun, 9:1}
\label{sec:unun-91}
\index{unun!9:1}
\index{9:1}
\index{transformator!9:1}

Tre lika lindningar kopplas i serie på samma kärna, trifilärt så att de
kopplar tätt. Skärmen och motvikten sitter i botten, mittledaren på
uttaget efter en tredjedel och antenntråden i toppen. Spänningen blir
tre gånger den koaxen lämnar, och impedansen nio gånger, alltså
50\,$\Omega$ mot omkring 450\,$\Omega$.

Det är anpassaren till den ändmatade tråd som med avsikt inte är
resonant, den som klarar sig med omkring 9:1 och ofta en tuner. En
halvvågs ändmatning ligger på några tusen ohm och är ett annat fall.
Utan motvikt blir radion själv den andra halvan av antennen.

   \begin{figure}[H]
    \centering
    \begin{tikzpicture}[
      x=1cm, y=1cm,
      ledareA/.style={rhbrav, line width=1.10pt, line cap=round, line join=round},
      ledareB/.style={rhbes, line width=1.10pt, line cap=round, line join=round},
      ledareC/.style={rhbtack, line width=1.10pt, line cap=round, line join=round},
      ledareAbak/.style={rhbrav!45, line width=0.85pt, line cap=round, dash pattern=on 2pt off 1.5pt},
      ledareBbak/.style={rhbes!45, line width=0.85pt, line cap=round, dash pattern=on 2pt off 1.5pt},
      ledareCbak/.style={rhbtack!45, line width=0.85pt, line cap=round, dash pattern=on 2pt off 1.5pt},
      koppling/.style={black!80, line width=0.8pt, line cap=round, line join=round},
      nod/.style={fill=black, circle, inner sep=1.3pt}
    ]

      % ====================================================================
      % VÄNSTER DEL: PRINCIPSCHEMA (TRE SERIEKOPPLADE LINDNINGAR, 1:3 SPÄNNING)
      % ====================================================================
      \node[etikett, font=\footnotesize\bfseries] at (-3.6, 3.20) {Principschema};

      % Ferritkärna (två vertikala linjer längs alla tre lindningarna)
      \draw[black!65, line width=1.1pt] (-3.15, -2.15) -- (-3.15, 2.45);
      \draw[black!65, line width=1.1pt] (-2.98, -2.15) -- (-2.98, 2.45);

      % Övre lindning A (blå, A1 -> A2)
      \draw[ledareA] (-3.35, 2.40)
        arc (90:-90:0.22 and 0.16)
        arc (90:-90:0.22 and 0.16)
        arc (90:-90:0.22 and 0.16)
        arc (90:-90:0.22 and 0.16);
      \fill[rhbrav] (-3.50, 2.27) circle[radius=0.05]; % Fasprick A1
      \node[etikett, text=rhbrav, anchor=east] at (-3.60, 2.27) {$A_1$};
      \node[etikett, text=rhbrav, anchor=east] at (-3.45, 1.25) {$A_2$};

      % Mellersta lindning B (rödbrun, B1 -> B2)
      \draw[ledareB] (-3.35, 0.86)
        arc (90:-90:0.22 and 0.16)
        arc (90:-90:0.22 and 0.16)
        arc (90:-90:0.22 and 0.16)
        arc (90:-90:0.22 and 0.16);
      \fill[rhbes] (-3.50, 0.73) circle[radius=0.05]; % Fasprick B1
      \node[etikett, text=rhbes, anchor=east] at (-3.60, 0.73) {$B_1$};
      \node[etikett, text=rhbes, anchor=east] at (-3.45, -0.30) {$B_2$};

      % Undre lindning C (grön, C1 -> C2)
      \draw[ledareC] (-3.35, -0.68)
        arc (90:-90:0.22 and 0.16)
        arc (90:-90:0.22 and 0.16)
        arc (90:-90:0.22 and 0.16)
        arc (90:-90:0.22 and 0.16);
      \fill[rhbtack] (-3.50, -0.81) circle[radius=0.05]; % Fasprick C1
      \node[etikett, text=rhbtack, anchor=east] at (-3.60, -0.81) {$C_1$};
      \node[etikett, text=rhbtack, anchor=east] at (-3.45, -1.85) {$C_2$};

      % Seriekoppling A2-B1 (2/3 upp) och B2-C1 (uttag 1/3 upp)
      \draw[koppling] (-3.35, 1.12) -- (-3.35, 0.86);
      \node[nod] (Uttag23) at (-3.35, 0.99) {};
      \node[etikett, anchor=west] at (-2.90, 0.99) {$A_2{+}B_1$};

      \draw[koppling] (-3.35, -0.42) -- (-3.35, -0.68);
      \node[nod] (Uttag13) at (-3.35, -0.55) {};

      % Koaxialkabel 50 ohm (nere till vänster, över lindning C)
      \fill[rhbjord!70, draw=black!70, line width=0.6pt] (-6.05, -1.55) rectangle (-4.95, -0.95);
      \draw[black!70, line width=0.6pt] (-4.95, -1.25) ellipse[x radius=0.08, y radius=0.30];
      \node[etikett, anchor=south] at (-5.50, -0.88) {Koax $50\,\Omega$};

      % Mittledare till uttaget efter 1/3 från botten (B2 + C1)
      \draw[koppling] (-4.95, -1.25) -- (-4.35, -1.25) -- (-4.35, -0.55) -- (Uttag13);
      \node[etikett, anchor=south] at (-4.25, -0.52) {Mittledare ($1/3$)};

      % Skärm + Motvikt i botten (C2)
      \draw[koppling] (-5.15, -1.55) -- (-5.15, -1.96) -- (-1.40, -1.96);
      \node[nod] at (-5.15, -1.55) {};
      \node[nod] at (-3.35, -1.96) {};
      \draw[black, fill=white, line width=0.7pt] (-1.40, -1.96) circle[radius=0.07];
      \node[etikett, anchor=north] at (-4.40, -2.00) {Skärm};
      \node[etikett, anchor=west] at (-1.25, -1.96) {Motvikt / jord};

      % Antenntråd i toppen (A1)
      \draw[koppling] (-3.35, 2.40) -- (-1.40, 2.40);
      \draw[black, fill=white, line width=0.7pt] (-1.40, 2.40) circle[radius=0.07];
      \node[etikett, anchor=west] at (-1.25, 2.40) {Antenntråd};

      % Impedansmarkering 450 ohm mellan antenntråd och motvikt
      \draw[{Stealth[length=1.6mm]}-{Stealth[length=1.6mm]}, black!60, line width=0.5pt]
        (-1.80, -1.80) -- (-1.80, 2.24)
        node[midway, fill=white, etikett] {$\approx 450\,\Omega$};


      % ====================================================================
      % HÖGER DEL: PRAKTISK TRIFILÄR LINDNING PÅ TOROIDKÄRNA
      % ====================================================================
      \begin{scope}[shift={(3.25, 0.65)}]
        \node[etikett, font=\footnotesize\bfseries] at (0, 2.55) {Trifilär lindning på toroidkärna};

        \def\Rout{1.85}
        \def\Rin{1.02}

        % Baksida av lindningsvarven (streckade under kärnan, 3 trådar per varv)
        \foreach \ang/\nextang in {195/157, 157/119, 119/81, 81/43, 43/5} {
          \draw[ledareAbak] ({\ang-2}:\Rin)  -- ({\nextang+4}:\Rout);
          \draw[ledareBbak] ({\ang-7}:\Rin)  -- ({\nextang-1}:\Rout);
          \draw[ledareCbak] ({\ang-12}:\Rin) -- ({\nextang-6}:\Rout);
        }

        % Toroidkärna (ring)
        \fill[rhbjord!75, even odd rule, draw=black!65, line width=0.8pt]
          (0,0) circle[radius=\Rout]
          (0,0) circle[radius=\Rin];

        % Framsida av lindningsvarven (tre trådar A, B och C parallellt)
        \foreach \ang in {195, 157, 119, 81, 43, 5} {
          \draw[ledareA] ({\ang+4}:\Rout) -- ({\ang-2}:\Rin);
          \draw[ledareB] ({\ang-1}:\Rout) -- ({\ang-7}:\Rin);
          \draw[ledareC] ({\ang-6}:\Rout) -- ({\ang-12}:\Rin);
        }

        % Startändar nere till vänster (från ytterkant: A1, B1, C1)
        \coordinate (StartA) at (199:\Rout);
        \coordinate (StartB) at (194:\Rout);
        \coordinate (StartC) at (189:\Rout);

        % Slutändar nere till höger (från innerkant: A2, B2, C2)
        \coordinate (SlutA) at ( 3:\Rin);
        \coordinate (SlutB) at (-2:\Rin);
        \coordinate (SlutC) at (-7:\Rin);

        % 1) A1 -> Antenntråd (uppe/vänster ut)
        \draw[ledareA] (StartA) -- (-2.35, -0.60) -- (-2.35, -1.35) -- (-3.05, -1.35);
        \draw[black, fill=white, line width=0.7pt] (-3.05, -1.35) circle[radius=0.07];
        \node[etikett, anchor=east] at (-3.15, -1.35) {Antenntråd};
        \node[etikett, text=rhbrav, anchor=east] at (-2.42, -0.80) {$A_1$};

        % 2) A2 + B1 -> Ihopkopplade i serie (inget yttre uttag)
        \coordinate (NodeAB) at (-0.45, -1.30);
        \draw[ledareB] (StartB) .. controls (-1.50, -0.95) and (-0.95, -1.30) .. (NodeAB);
        \draw[ledareA] (SlutA)  .. controls ( 0.75, -0.45) and ( 0.15, -1.30) .. (NodeAB);
        \node[nod] at (NodeAB) {};
        \node[etikett, anchor=south] at (-0.45, -1.25) {Serie ($A_2{+}B_1$)};

        % 3) B2 + C1 -> Koax mittledare (uttag efter 1/3)
        \coordinate (NodeBC) at (-0.55, -1.85);
        \draw[ledareC] (StartC) .. controls (-1.45, -1.45) and (-0.95, -1.85) .. (NodeBC);
        \draw[ledareB] (SlutB)  .. controls ( 0.85, -0.85) and ( 0.10, -1.85) .. (NodeBC);
        \node[nod] at (NodeBC) {};
        \draw[koppling] (NodeBC) -- (-0.55, -2.55);
        \node[etikett, text=rhbes, anchor=east] at (-0.65, -2.10) {$B_2{+}C_1$};

        % 4) C2 -> Koax skärm + Motvikt (i botten)
        \coordinate (NodeC2) at (0.15, -2.15);
        \draw[ledareC] (SlutC) .. controls (1.35, -0.55) and (1.15, -2.15) .. (NodeC2);
        \node[nod] at (NodeC2) {};
        \draw[koppling] (NodeC2) -- (0.15, -2.55);
        \draw[koppling] (NodeC2) -- (2.75, -2.15);
        \draw[black, fill=white, line width=0.7pt] (2.75, -2.15) circle[radius=0.07];
        \node[etikett, anchor=west] at (2.85, -2.15) {Motvikt};
        \node[etikett, text=rhbtack, anchor=south west] at (1.15, -1.45) {$C_2$};

        % Koaxsymbol nedtill
        \fill[rhbjord!70, draw=black!70, line width=0.6pt] (-0.80, -3.15) rectangle (0.40, -2.55);
        \node[nod] at (0.15, -2.55) {};
        \node[etikett, anchor=east] at (-0.88, -2.85) {Mittledare ($B_2{+}C_1$)};
        \node[etikett, anchor=west] at ( 0.48, -2.85) {Skärm ($C_2$)};
        \node[etikett, anchor=north] at (-0.20, -3.20) {Koax $50\,\Omega$};
      \end{scope}

    \end{tikzpicture}
    \caption{Unun 9:1 för icke-resonant ändmatad tråd (EFRW). Tre lika trådar
    ($A$, $B$ och $C$) lindas trifilärt åt samma håll och seriekopplas ($A_2$ med
    $B_1$, samt $B_2$ med $C_1$). Skärm och motvikt läggs i botten på $C_2$,
    koaxens mittledare på uttaget efter en tredjedel ($B_2{+}C_1$) och
    antenntråden i toppen på $A_1$.}
    \label{fig:unun-91}
    \end{figure}


\subsubsection{Unun, 49:1}
\label{sec:unun-491}
\index{unun!49:1}
\index{49:1}
\index{transformator!49:1}
\index{endmatad!49:1}

Här är det en enda lindning med ett uttag, inte tre lika delar.
Antenntråden ser sju gånger så många varv som den del radion ligger
över. Sju i kvadrat är 49, så några tusen ohm, omkring 2\,500\,$\Omega$,
kommer ner mot 50\,$\Omega$. I praktiken ligger uttaget ofta vid två
varv när hela lindningen är fjorton, vilket är samma sju gånger.

Det är transformatorn till den ändmatade halvvågen. Förlusten kommer
av att den höga sidan är många varv, och tråden kan inte vara grov om
den ska få plats och ändå koppla. En motvikt hör till, eller koaxens
skärm om en strömbalun 1:1 sitter en bit ner på kabeln, så att skärmen
inte leder common mode-strömmen hela vägen in i rummet. På 15 och
10\,m lägger många en kondensator på omkring 100\,pF över
50\,$\Omega$-sidan, eftersom uttaget annars lämnar de banden högt i
stående våg.

    \begin{figure}[H]
    \centering
    \begin{tikzpicture}[
      x=1cm, y=1cm,
      ledareA/.style={rhbrav, line width=1.15pt, line cap=round, line join=round},
      ledareB/.style={rhbes, line width=1.15pt, line cap=round, line join=round},
      ledareAbak/.style={rhbrav!45, line width=0.85pt, line cap=round, dash pattern=on 2pt off 1.5pt},
      ledareBbak/.style={rhbes!45, line width=0.85pt, line cap=round, dash pattern=on 2pt off 1.5pt},
      koppling/.style={black!80, line width=0.8pt, line cap=round, line join=round},
      nod/.style={fill=black, circle, inner sep=1.3pt}
    ]

      % ====================================================================
      % VÄNSTER DEL: PRINCIPSCHEMA (2 + 12 = 14 VARV, OMSÄTTNING 1:49)
      % ====================================================================
      \node[etikett, font=\footnotesize\bfseries] at (-3.6, 3.20) {Principschema};

      % Ferritkärna (två vertikala linjer)
      \draw[black!65, line width=1.1pt] (-3.15, -2.05) -- (-3.15, 2.45);
      \draw[black!65, line width=1.1pt] (-2.98, -2.05) -- (-2.98, 2.45);

      % Övre lindningsdel (12 varv, rödbrun, från uttaget vid 2 varv upp till varv 14)
      \draw[ledareB] (-3.35, 2.40)
        arc (90:-90:0.23 and 0.25)
        arc (90:-90:0.23 and 0.25)
        arc (90:-90:0.23 and 0.25)
        arc (90:-90:0.23 and 0.25)
        arc (90:-90:0.23 and 0.25)
        arc (90:-90:0.23 and 0.25);
      \node[etikett, text=rhbes, anchor=east] at (-3.55, 0.90) {$12$ varv};
      \node[etikett, text=rhbes, anchor=east] at (-3.45, 2.25) {varv 14};

      % Undre lindningsdel (2 varv, blå, från botten/jord upp till uttaget vid 2 varv)
      \draw[ledareA] (-3.35, -0.86)
        arc (90:-90:0.23 and 0.27)
        arc (90:-90:0.23 and 0.27);
      \node[etikett, text=rhbrav, anchor=east] at (-3.55, -1.40) {$2$ varv};

      % Uttag vid 2 varv
      \draw[koppling] (-3.35, -0.60) -- (-3.35, -0.86);
      \node[nod] (Uttag2) at (-3.35, -0.73) {};

      % Koaxialkabel 50 ohm (nere till vänster)
      \fill[rhbjord!70, draw=black!70, line width=0.6pt] (-6.30, -1.65) rectangle (-5.25, -1.05);
      \draw[black!70, line width=0.6pt] (-5.25, -1.35) ellipse[x radius=0.08, y radius=0.30];
      \node[etikett, anchor=south] at (-5.78, -0.98) {Koax $50\,\Omega$};

      % Mittledare till uttaget vid 2 varv
      \draw[koppling] (-5.25, -1.35) -- (-4.75, -1.35) -- (-4.75, -0.73) -- (Uttag2);
      \node[etikett, anchor=south] at (-4.15, -0.70) {Uttag (varv 2)};

      % Skärm + Motvikt i botten (varv 0)
      \draw[koppling] (-5.45, -1.65) -- (-5.45, -1.94) -- (-1.40, -1.94);
      \node[nod] at (-5.45, -1.65) {};
      \node[nod] at (-3.35, -1.94) {};
      \draw[black, fill=white, line width=0.7pt] (-1.40, -1.94) circle[radius=0.07];
      \node[etikett, anchor=north] at (-4.85, -1.98) {Skärm};
      \node[etikett, anchor=west] at (-1.25, -1.94) {Motvikt / jord};

      % Kompensationskondensator 100 pF över 50-ohmssidan
      \draw[koppling] (-4.35, -0.73) -- (-4.35, -1.24);
      \draw[koppling] (-4.35, -1.42) -- (-4.35, -1.94);
      \draw[black!85, line width=1.0pt] (-4.55, -1.24) -- (-4.15, -1.24);
      \draw[black!85, line width=1.0pt] (-4.55, -1.42) -- (-4.15, -1.42);
      \node[nod] at (-4.35, -0.73) {};
      \node[nod] at (-4.35, -1.94) {};
      \node[etikett, anchor=east] at (-4.58, -1.58) {$100\,\text{pF}$};

      % Antenntråd (EFHW) i toppen (varv 14)
      \draw[koppling] (-3.35, 2.40) -- (-1.40, 2.40);
      \draw[black, fill=white, line width=0.7pt] (-1.40, 2.40) circle[radius=0.07];
      \node[etikett, anchor=west] at (-1.25, 2.40) {Antenntråd ($\lambda/2$)};

      % Impedansmarkering 2450 ohm mellan antenntråd och motvikt
      \draw[{Stealth[length=1.6mm]}-{Stealth[length=1.6mm]}, black!60, line width=0.5pt]
        (-1.80, -1.78) -- (-1.80, 2.24)
        node[midway, fill=white, etikett] {$\approx 2\,450\,\Omega$};


      % ====================================================================
      % HÖGER DEL: PRAKTISK LINDNING PÅ TOROIDKÄRNA (2 BIFILÄRA + 12 ENKLA)
      % ====================================================================
      \begin{scope}[shift={(3.25, 0.55)}]
        \node[etikett, font=\footnotesize\bfseries] at (0, 2.65) {Lindning på toroidkärna ($2 + 12 = 14$ varv)};

        \def\Rout{1.85}
        \def\Rin{1.02}

        % Baksida av de första 2 bifilära varven (nere till vänster)
        \foreach \ang/\nextang in {224/204, 204/184} {
          \draw[ledareAbak] ({\ang-2}:\Rin) -- ({\nextang+2}:\Rout);
          \draw[ledareBbak] ({\ang+3}:\Rin) -- ({\nextang+7}:\Rout);
        }

        % Baksida av de återstående 12 enkelvarven runt kärnan
        \foreach \ang/\nextang in {184/164, 164/144, 144/124, 124/104, 104/84, 84/64, 64/44, 44/24, 24/4, 4/-16, -16/-36} {
          \draw[ledareBbak] (\ang:\Rin) -- (\nextang:\Rout);
        }

        % Toroidkärna (ring)
        \fill[rhbjord!75, even odd rule, draw=black!65, line width=0.8pt]
          (0,0) circle[radius=\Rout]
          (0,0) circle[radius=\Rin];

        % Framsida av de första 2 bifilära varven (blå + rödbrun tätt ihop)
        \foreach \ang in {224, 204} {
          \draw[ledareA] ({\ang+2}:\Rout) -- ({\ang-2}:\Rin);
          \draw[ledareB] ({\ang+7}:\Rout) -- ({\ang+3}:\Rin);
        }

        % Framsida av de 12 återstående varven (enkel rödbrun tråd varv 3..14)
        \foreach \ang in {184, 164, 144, 124, 104, 84, 64, 44, 24, 4, -16, -36} {
          \draw[ledareB] (\ang:\Rout) -- (\ang:\Rin);
        }

        % Start i botten (varv 0: både blå och rödbrun börjar i jord/skärm)
        \coordinate (StartA) at (226:\Rout);
        \coordinate (StartB) at (231:\Rout);

        % Uttag efter 2 varv (blå tråd slutar vid varv 2 och kopplas till mittledaren)
        \coordinate (TapA)   at (184:\Rout);

        % Slut efter 14 varv (nere till höger -> Antenntråd)
        \coordinate (SlutB)  at (-36:\Rin);

        % 1) Uttag efter 2 varv -> Koax mittledare
        \draw[ledareA] (TapA) -- (-2.35, -0.13) -- (-2.35, -2.05) -- (-0.55, -2.05) -- (-0.55, -2.50);
        \node[nod] at (TapA) {};
        \node[etikett, text=rhbrav, anchor=east] at (-2.42, -0.35) {Uttag (2 varv)};

        % 2) Start i botten (varv 0) -> Koax skärm + Motvikt
        \coordinate (JordPunkt) at (0.10, -2.05);
        \draw[ledareA] (StartA) .. controls (-1.35, -1.75) and (-0.45, -2.05) .. (JordPunkt);
        \draw[ledareB] (StartB) .. controls (-1.15, -1.65) and (-0.35, -2.05) .. (JordPunkt);
        \node[nod] at (JordPunkt) {};
        \draw[koppling] (JordPunkt) -- (0.10, -2.50);
        \draw[koppling] (JordPunkt) -- (1.85, -2.05);
        \draw[black, fill=white, line width=0.7pt] (1.85, -2.05) circle[radius=0.07];
        \node[etikett, anchor=west] at (1.95, -2.05) {Motvikt};

        % 3) Kondensator 100 pF mellan mittledare och skärm
        \draw[koppling] (-1.45, -2.05) -- (-1.45, -2.35);
        \draw[black!85, line width=0.95pt] (-1.62, -2.35) -- (-1.28, -2.35);
        \draw[black!85, line width=0.95pt] (-1.62, -2.48) -- (-1.28, -2.48);
        \draw[koppling] (-1.45, -2.48) -- (-1.45, -2.72) -- (0.10, -2.72);
        \node[nod] at (-1.45, -2.05) {};
        \node[nod] at (0.10, -2.72) {};
        \node[etikett, anchor=east] at (-1.68, -2.42) {$100\,\text{pF}$};

        % 4) Slut efter 14 varv -> Antenntråd (EFHW)
        \draw[ledareB] (SlutB) .. controls (1.35, -0.75) and (1.95, -1.05) .. (2.25, -1.30) -- (2.85, -1.30);
        \draw[black, fill=white, line width=0.7pt] (2.85, -1.30) circle[radius=0.07];
        \node[etikett, anchor=west] at (2.95, -1.30) {Antenntråd};
        \node[etikett, text=rhbes, anchor=south west] at (1.55, -0.92) {varv 14};
        \node[etikett, text=rhbes] at (0, 0) {$2 + 12$ varv};

        % Koaxsymbol nedtill
        \fill[rhbjord!70, draw=black!70, line width=0.6pt] (-0.80, -3.15) rectangle (0.35, -2.50);
        \node[nod] at (0.10, -2.50) {};
        \node[etikett, anchor=east] at (-0.88, -2.85) {Mittledare (varv 2)};
        \node[etikett, anchor=west] at ( 0.43, -2.85) {Skärm (varv 0)};
        \node[etikett, anchor=north] at (-0.22, -3.20) {Koax $50\,\Omega$};
      \end{scope}

    \end{tikzpicture}
    \caption{Unun 49:1 för ändmatad halvvåg (EFHW). De två första varven lindas
    bifilärt från jord/skärm upp till uttaget för koaxens mittledare vid varv~2,
    varefter antenntrådens lindning fortsätter ytterligare 12~varv (totalt
    14~varv, spänningsomsättning $1:7$ och impedansomsättning $1:49$). En
    kondensator på $100\,\text{pF}$ över ingången kompenserar reaktansen på de
    högre banden.}
    \label{fig:unun-491}
    \end{figure}


